% !TeX program = xelatex
% ============================================================
%  OpenAI 四维挂谷猜想工作综述
%  数学思想脉络、核心技术谱系与验证状态
%  作者：TRAE-Work、Deepseek-V4.1-Flash
%  编译：xelatex 多遍自动收敛（见 build.ps1）
% ============================================================
\documentclass[11pt,a4paper]{ctexart}

\usepackage{geometry}
\geometry{left=2.4cm,right=2.4cm,top=2.6cm,bottom=2.6cm}
\usepackage{amsmath,amssymb,amsthm}
\usepackage{array,booktabs,tabularx,multirow}
\usepackage[dvipsnames,table]{xcolor}
\usepackage{enumitem}
\usepackage[numbers,sort&compress]{natbib}
\usepackage[most]{tcolorbox}
\usepackage{tikz}
\usetikzlibrary{arrows.meta,positioning,fit,calc}
\usepackage{caption}
\usepackage[section]{placeins}
\usepackage{hyperref}

% ---------- 配色（藏青主色 + 天蓝辅助 + 浅底） ----------
\definecolor{navy}{HTML}{1F3864}
\definecolor{sky}{HTML}{2E75B6}
\definecolor{lightbg}{HTML}{F2F6FB}
\definecolor{lightline}{HTML}{D6E4F0}
\definecolor{warnc}{HTML}{9C4221}
\definecolor{tealc}{HTML}{1F6F63}
\definecolor{greyx}{HTML}{5A6472}

% ---------- 中英文间距 0.22em ----------
\xeCJKsetup{CJKecglue={\hskip 0.22em}}

% ---------- 版式 ----------
\linespread{1.35}
\setlength{\parindent}{2em}
\renewcommand{\arraystretch}{1.35}

% 浮动体参数：允许每页多放表格，避免堆积到文末
\setcounter{topnumber}{3}
\setcounter{bottomnumber}{2}
\setcounter{totalnumber}{4}
\renewcommand{\topfraction}{0.92}
\renewcommand{\bottomfraction}{0.5}
\renewcommand{\textfraction}{0.08}
\renewcommand{\floatpagefraction}{0.75}

% 表格列类型：raggedright，避免窄列单元格拉伸产生 Underfull
\newcolumntype{L}[1]{>{\raggedright\arraybackslash}p{#1}}
\newcolumntype{Y}{>{\raggedright\arraybackslash}X}

\ctexset{
  section = {format = {\color{navy}\bfseries\large}, beforeskip = {14pt plus 3pt minus 2pt}, afterskip = {9pt plus 2pt minus 2pt}},
  subsection = {format = {\color{navy}\bfseries}, beforeskip = {10pt plus 2pt minus 2pt}, afterskip = {6pt plus 1pt minus 1pt}},
  subsubsection = {format = {\color{sky}\bfseries}, beforeskip = {8pt plus 2pt minus 2pt}, afterskip = {5pt plus 1pt minus 1pt}}
}

% ---------- 定理环境（标题上色，正文正体） ----------
\newtheoremstyle{navystyle}{0.6em}{0.6em}{\normalfont}{}{\color{navy}\bfseries}{}{0.6em}{}
\newtheoremstyle{skystyle}{0.6em}{0.6em}{\normalfont}{}{\color{sky}\bfseries}{}{0.6em}{}
\theoremstyle{navystyle}
\newtheorem{theorem}{定理}[section]
\newtheorem{proposition}[theorem]{命题}
\newtheorem{lemma}[theorem]{引理}
\newtheorem{corollary}[theorem]{推论}
\theoremstyle{skystyle}
\newtheorem{definition}[theorem]{定义}
\newtheorem{remark}[theorem]{注}
\newtheorem{example}[theorem]{例}

% ---------- 卡片（标签独立成行，避免与正文重叠） ----------
\newtcolorbox{keybox}[1][关键判断]{breakable, enhanced, frame hidden, colback = lightbg,
  left = 3.2mm, right = 2.4mm, top = 0.8mm, bottom = 1.6mm,
  borderline west = {1.6pt}{0pt}{navy},
  coltitle = navy, colbacktitle = lightbg, fonttitle = \bfseries\small,
  titlerule = 0pt, title = {#1},
  before skip = 7pt, after skip = 7pt,
  fontupper = \normalfont\small
}
\newtcolorbox{evbox}[1][证据]{breakable, enhanced, frame hidden, colback = lightbg,
  left = 3.2mm, right = 2.4mm, top = 0.8mm, bottom = 1.6mm,
  borderline west = {1.6pt}{0pt}{sky},
  coltitle = sky, colbacktitle = lightbg, fonttitle = \bfseries\small,
  titlerule = 0pt, title = {#1},
  before skip = 7pt, after skip = 7pt,
  fontupper = \normalfont\small
}
\newtcolorbox{cautbox}[1][边界与不确定性]{breakable, enhanced, frame hidden, colback = lightbg,
  left = 3.2mm, right = 2.4mm, top = 0.8mm, bottom = 1.6mm,
  borderline west = {1.6pt}{0pt}{warnc},
  coltitle = warnc, colbacktitle = lightbg, fonttitle = \bfseries\small,
  titlerule = 0pt, title = {#1},
  before skip = 7pt, after skip = 7pt,
  fontupper = \normalfont\small
}
\newtcolorbox{dirbox}[1][方向建议]{breakable, enhanced, frame hidden, colback = lightbg,
  left = 3.2mm, right = 2.4mm, top = 0.8mm, bottom = 1.6mm,
  borderline west = {1.6pt}{0pt}{tealc},
  coltitle = tealc, colbacktitle = lightbg, fonttitle = \bfseries\small,
  titlerule = 0pt, title = {#1},
  before skip = 7pt, after skip = 7pt,
  fontupper = \normalfont\small
}

% ---------- 置信度标签 ----------
\newcommand{\cfH}{\textcolor{navy}{\textbf{〔置信度：高〕}}}
\newcommand{\cfM}{\textcolor{sky}{\textbf{〔置信度：中〕}}}
\newcommand{\cfL}{\textcolor{warnc}{\textbf{〔置信度：低〕}}}

% ---------- 上标引用 ----------
\newcommand{\upcite}[1]{\textsuperscript{\citep{#1}}}
\newcommand{\upcites}[1]{\textsuperscript{[#1]}}

% ---------- 参考文献 ----------
\renewcommand{\refname}{参考文献}
\newcommand{\bibtype}[1]{[#1]}

% ---------- 超链接 ----------
\hypersetup{
  colorlinks = true,
  linkcolor = navy,
  citecolor = sky,
  urlcolor = sky,
  bookmarksnumbered = true,
  pdftitle = {OpenAI 四维挂谷猜想工作综述：数学思想脉络、核心技术谱系与验证状态},
  pdfauthor = {TRAE-Work, Deepseek-V4.1-Flash}
}

\begin{document}

% ---------- 标题区 ----------
\begin{center}
  {\LARGE\bfseries\color{navy} OpenAI 四维挂谷猜想工作综述}\\[7pt]
  {\large\color{navy} 数学思想脉络、核心技术谱系与验证状态}\\[12pt]
  {\normalsize TRAE-Work、Deepseek-V4.1-Flash}\\[12pt]
  {\color{navy}\rule{0.86\textwidth}{1.1pt}}\\[2pt]
  {\color{sky}\rule{0.86\textwidth}{0.5pt}}
\end{center}

\vspace{4pt}
\noindent\textbf{\color{navy}摘\quad 要}：本文以 OpenAI 于 2026 年 9 月 23 日、24 日发布的两篇挂谷猜选手稿为主要对象：一篇声称证明三维挂谷极大函数猜想，把 Wang--Zahl 定理留下的密度幂缺口 $\lambda^{K(\varepsilon)}$ 补到尖锐的 $\lambda^{3}$；另一篇声称证明四维挂谷集的 Hausdorff 维数等于 $4$，且不设可测性、紧性与粘性假设。两篇手稿均未经同行评审，也未纳入 Lean 形式化清单，仓库自述未形式化的结果可能存在问题，因此本文的第一项工作是事实核验：整理发布时间线、稿件参数、依赖链与验证状态，并指出四维结果的关键输入取自三维姊妹篇的一条加权板条引理，其有效性无法与三维篇的核验状态切割开来。第二项工作是技术拆解：从任意挂谷集到加权模型的归约、以“地平线”为载体的时间代价记账、带质量约束的二次图卡、极值临界路径与三窄度区间的分类讨论，逐层展开四维证明的机理，并说明它与三维范式（粘性、平面性、颗粒性）在对象层面的区别。第三项工作是技术谱系分析：把核心技术分为直系继承、门内改造、跨门派借入与谱系内首次系统使用四类，逐条给出证据与置信度。对照 Wang--Zahl 与 Guth--Wang--Zahl 原文可以确认，手稿使用的熵账本、概率耦合、实代数几何的定量纪律与乘法卷积的 Fourier 衰减估计，在传统挂谷工具箱中并无先例。最后给出可复现的核验路线图，以及验证、数学推进、方法回灌、工程流程四个方向的下一步建议。全文对手稿声称与已证定理严格区分。
\vspace{4pt}

\noindent\textbf{\color{navy}关键词}：挂谷猜想；Kakeya 集；Hausdorff 维数；极大函数；调和分析；AI 辅助数学；形式化验证

\vspace{4pt}
\noindent\textbf{\color{navy}MSC (2020)}：42B25；28A78；42B10；68V20

\vspace{6pt}
\noindent\textbf{\color{navy}Abstract.}\ This survey examines the two Kakeya manuscripts released by OpenAI on September 23 and 24, 2026. The first claims to prove the three-dimensional Kakeya maximal conjecture, sharpening the density exponent $\lambda^{K(\varepsilon)}$ left open after Wang--Zahl to the critical $\lambda^{3}$; the second claims that every four-dimensional Kakeya set has Hausdorff dimension four, with no measurability, compactness, or stickiness hypothesis. Neither manuscript has been peer-reviewed or formalized in Lean, and the repository states that unformalized results may contain issues. The survey proceeds in four parts. First, a verification file: timeline, manuscript parameters, dependency chain, and verification status, including the observation that the four-dimensional result imports a weighted plank bound from its three-dimensional companion, so its validity is coupled to that of the companion. Second, a technical reconstruction: the reduction from an arbitrary Kakeya set to a weighted model, the bookkeeping of a positive ``horizon'' (a time-cost obstruction), quadratic charts with mass constraints, extremal critical paths, and the three narrowness regimes. Third, a provenance analysis that classifies each core technique as direct inheritance, in-school modification, cross-discipline import, or first systematic use within the Kakeya lineage; comparison with Wang--Zahl and Guth--Wang--Zahl shows that the entropy ledger, probabilistic coupling, quantitative real algebraic geometry, and Fourier-decay estimates have no precedent in the traditional Kakeya toolkit. Fourth, a reproducible verification roadmap and directions for future work along four tracks: verification, mathematics, methodology transfer, and research workflow. Claims of the manuscripts are kept strictly separate from established theorems, and each judgement carries a stated confidence level.

\vspace{4pt}
\noindent\textbf{\color{navy}Keywords}：Kakeya conjecture; Kakeya set; Hausdorff dimension; maximal function; harmonic analysis; AI for mathematics; formal verification

\vspace{10pt}
\noindent\textbf{\color{navy}阅读约定}：本文所有涉及 OpenAI 稿件的结论一律表述为“手稿声称”，与已证定理严格区分；涉及事实性断言的句子以上标编号给出出处；分析性判断在正文中明示“本文判断”字样。数学记号见附录 A。

% ---------- 目录 ----------
\vspace{8pt}
{\setlength{\parskip}{2pt}\tableofcontents}

\clearpage

\section{导言：问题的坐标与本文的任务}

2026 年 10 月 6 日，OpenAI 在一个公开仓库中放出了内部前沿模型生成的数学手稿合集，共 372 个结果族；翌日的修订记录撤回了其中 3 篇并修订 14 篇，其余保留，此后按版本号管理修正\upcite{openairelease2026,openairepo2026}。这批材料里，第 074 族的体量不算最大，位置却最扎眼：它只有两篇稿件，主题是挂谷（Kakeya）猜想。一篇 97 页，声称证明三维挂谷极大函数猜想；另一篇 175 页，声称证明四维挂谷集的 Hausdorff 维数等于 $4$\upcite{openai3d2026,openai4d2026}。作者栏里没有个人署名，只有一个机构名。

值得为这一族单独写一份综述，理由不只是“AI 又解了一个大问题”。挂谷猜想的三维情形在 2025 年 2 月刚被王虹与 Zahl 解决，2026 年 7 月费城国际数学家大会上，这项工作成为王虹获菲尔兹奖的主要成果之一\upcite{wangz2025,xinhua2026,pku2026}。这是一条仍在高速推进、社区记忆新鲜的研究线，OpenAI 手稿直接踩在这条线的最前沿，而不是挑选一个冷门角落。另一方面，两篇稿件都没有进入仓库的 Lean 形式化清单，也没有经过同行评审；仓库自己的措辞是“部分未形式化的结果可能存在问题”\upcite{openairepo2026}。结果的真伪悬置，方法的面貌却已经摊开。这正是综述可以介入的时机。

本文的任务有六项。其一，把发布事实、稿件参数、依赖链条与验证状态整理成一份可核对的档案（第 4 章），其中最关键的一条是：四维稿件的核心输入取自三维姊妹篇的一条加权板条引理，两篇稿件在验证意义上是一体的，不能拆开评估。其二，逐层重建四维证明的技术骨架（第 5 章），覆盖从任意挂谷集到加权模型的归约、“地平线”机制、图卡体系、极值临界路径与三窄度区间的分类讨论。其三，做技术谱系分析（第 6、7 章）：把全部核心技术按直系继承、门内改造、跨门派借入、谱系内首次系统使用四类逐一归档，并给出证据与置信度。其四，把这条研究线上的主要方法放进一张统一比较矩阵（第 8 章）。其五，分析验证缺口（第 9 章），说明“未经形式化、思想未经独立推演”对这份工作的具体含义。其六，给出核验路线图与下一步方向建议（第 10 章），并在收束处给出修正后的判断（第 11 章）。

本文不做三件事。不做 175 页的逐行核验——那需要一支熟悉四维管族几何的验证小组，超出单篇综述的职责与能力；不复述稿件中的全部定义与引理编号，只保留理解机制所必需的骨架；不代替数学界裁决正确性，所有涉及稿件的结论一律写作“手稿声称”。

材料与方法是本文可信度的全部来源，因此一并说明。一手材料包括两篇稿件 PDF 全文、仓库的目录清单、Lean 形式化清单与修订史、OpenAI 官方发布说明\upcite{openai3d2026,openai4d2026,openairepo2026,openairelease2026}；历史与工具箱基线取自 Zahl 的挂谷综述（覆盖 2000--2025 年）\upcite{zahl2025survey}；为判断“哪些工具是门内已用、哪些是外门派借入”，本文对 Wang--Zahl 的三维论文与 Guth--Wang--Zahl 的简化证明两份门内主力文献做了全文关键词扫描，结果在第 7 章给出。事实性断言以上标编号标注出处；属于本文分析判断的，在文中明示，并按〔置信度：高/中/低〕标注。技术性转述以稿件自身的章节与命题编号为准（如 Theorem 2.17、Proposition 6.8），读者可以直接按编号定位原文核对，这也用于防止转述层引入不必要的失真。

导言的最后一段话说明本文的立场。挂谷猜想在四维是否已经被解决，本文不给出结论；本文给出的是这份稿件的结构、来源、缺口与可能去处。一份正确但无人理解的证明，和一份无人能验证的声称，在数学共同体里的地位完全不同；把二者区分开，是接下来十章要做的事。
\section{本质问题：方向完备性与空间压缩}

\begin{keybox}[本质问题]
挂谷型问题的全部内容，是“方向完备性”与“空间压缩”这两条要求之间的张力：一个集合必须在每个方向上装得下一根单位线段（方向完备），却又不许它因此获得正测度（空间压缩的极限）。测度被压到零之后，能继续度量“大小”的只剩下维数，于是问题转化为：方向完备的性能否把维数钉在空间维数 $n$ 上。
\end{keybox}

\subsection{从针的翻转问题到方向完备集合}

挂谷问题有两个源头，通常被合在一起叙述，实际是两回事。1917 年 Fujiwara 与挂谷宗一研究了定宽曲线的最大值最小值问题，其中包含一个几何问句：在平面上让一根单位长的小针连续转动、两端点互换位置，扫过的区域面积最小是多少\upcite{fujiwara1917}。Besicovitch 在 1928 年给出反直觉的答案：这样的区域面积可以任意小，且存在包含每个方向线段、但 Lebesgue 测度为零的集合\upcite{besicovitch1928}。这是两个不同的陈述——前者关于连续运动，后者关于“每个方向装得下一根线段”这种静态的方向完备性；二者经常被混为一谈，却由不同的构造支撑。

现代问题取的是后一种表述：$\mathbb{R}^n$ 中的一个挂谷集（Kakeya set，又称 Besicovitch 集）是包含每个方向一条单位线段的集合。Besicovitch 的构造说明挂谷集可以测度为零，于是“面积”这把尺子失效。

\subsection{测度失效之后，维数成为新的尺子}

对零测集合，“大小”要用维数来度量，而维数并非一种尺度的简单替换，它描述的是集合在不同放大倍率下的覆盖复杂度。$n$ 维空间中的一个有界集合 $E$，若用边长为 $\delta$ 的小立方体去覆盖，覆盖数 $N_\delta(E)$ 的行为刻画了 Minkowski 维数：
\begin{equation}
\underline{\dim}_{\mathrm{M}} E \;=\; \liminf_{\delta\to 0}\frac{\log N_\delta(E)}{\log (1/\delta)},\qquad
\overline{\dim}_{\mathrm{M}} E \;=\; \limsup_{\delta\to 0}\frac{\log N_\delta(E)}{\log (1/\delta)}.
\end{equation}
Minkowski 维数只在一个公共尺度上计数；Hausdorff 维数则允许覆盖件直径各不相同，更能察觉跨尺度的精细结构。对挂谷问题，Hausdorff 版本通常比 Minkowski 版本更难，因为“每个方向”的要求必须对任意尺度、任意覆盖都成立。

\begin{keybox}[第一性原理]
挂谷猜想断言的是：方向完备性足以把 Hausdorff 维数钉在满维 $n$。为什么这条断言不平凡？因为“包含每个方向一条线段”与“集合占满维数”之间并没有逻辑蕴含——平面上的直线族可以沿着低维代数曲面大量聚集，把体积压到零的同时不损失方向覆盖。真正阻止压缩的，不是测度，而是“重叠后必须仍然在每个方向上保持足够复杂度”的几何约束。
\end{keybox}

\subsection{把集合问题离散化：细管与阴影}

研究挂谷问题几乎总是先做离散化。固定小尺度 $\delta>0$，把单位线段加粗成半径为 $\delta$ 的细管（tube），每个 $\delta$-分离的方向取一根细管，得到一族约 $\delta^{1-n}$ 根细管。挂谷集猜想要求这些细管的并集在体积上达到临界下界。若再给每根细管指定一个可测子集（阴影），并要求阴影占细管体积的比例至少为 $\lambda$，就得到更强、也更具分析味道的极大函数版本：
\begin{equation}
\Bigl|\bigcup_{T\in\mathcal{T}} Y(T)\Bigr| \;\gtrsim_\varepsilon\; \delta^{\varepsilon}\, \lambda^{n}\sum_{T\in\mathcal{T}}|T|,
\end{equation}
其中 $Y(T)\subseteq T$ 为阴影。当 $\lambda$ 是常数量级时，上式给出维数结论；当 $\lambda$ 可以任意小、且指数精确为 $n$ 时，上式对应极大函数猜想。后者比前者更细：它要求对任意“填充比例”都保持同一幂次。

这个离散化视角还说明为什么挂谷问题与 Fourier 分析纠缠在一起。Fefferman 的球乘子反例显示，球乘子的 $L^p$ 无界性本质上来自管族在某些区域的高度重叠\upcite{fefferman1971}；Bourgain 进一步把极大函数估计与 Fourier 限制猜想联系起来\upcite{bourgain1991}。管族重叠行为的几何控制，成为限制理论、Bochner--Riesz 求和与波动方程局部光滑问题的公共枢纽。

\subsection{为什么四维会成为一道坎}

二维情形的集合版本由 Davies 解决\upcite{davies1971}，极大函数版本由 C\'ordoba 解决\upcite{cordoba1977}。三维在 2025 年被王虹--Zahl 拿下\upcite{wangz2025}。四维则不同：Guth 指出，三维证明中的关键定理搬到四维会失效，因为在更高维里，细管可以沿着二次曲面（如 $ad-bc=1$ 型超曲面）大量聚集，形成三维里不存在的危险结构\upcite{guthzahl2018,zahl2025survey}。这解释了为什么四维的维数下界长期停留在 3.059 附近（见第 3 章演进表），也解释了 OpenAI 四维稿件为何把“二次图卡局部拟合”作为主轴：它要正面回应“低次代数曲面上的聚集”这一四维特有障碍。

\begin{keybox}[可记忆主句]
测度可以消失，方向不可消失；维数，是方向完备性对空间压缩的最后一道防线。
\end{keybox}
\section{历史脉络与门派图谱：从挂谷宗一到王虹--Zahl}

四维稿件的每一件工具都能在历史上找到来处，而每一处来处又都对应一个研究门派。本章先把主干时间线铺一遍（表 \ref{tab:timeline}），再把门派与工具对号入座（表 \ref{tab:schools}），最后交代四维下界的演进位置（表 \ref{tab:r4bounds}）与一条贯穿全史的现象：挂谷问题本身就是一部跨门派借入史。

\subsection{主干时间线}

\begin{table}[t]
\centering
\caption{挂谷问题主干里程碑（1917--2026）}
\label{tab:timeline}
\footnotesize
\begin{tabularx}{\textwidth}{@{}L{1.5cm}L{4.2cm}Y@{}}
\toprule
\rowcolor{lightbg}\textbf{年份} & \textbf{进展} & \textbf{方法要点与出处} \\
\midrule
1917 & 针的翻转问题 & Fujiwara--挂谷宗一提出定宽曲线与针转动的最小面积问题\upcite{fujiwara1917} \\
1928 & 零测度挂谷集 & Besicovitch 构造面积任意小、可含全部方向的集合\upcite{besicovitch1928} \\
1971 & 平面情形解决 & Davies 证明平面挂谷集 Hausdorff 维数为 2\upcite{davies1971} \\
1971 & 与 Fourier 分析接通 & Fefferman 用管族重叠构造球乘子反例\upcite{fefferman1971} \\
1977 & 平面极大函数 & C\'ordoba 给出平面极大估计\upcite{cordoba1977} \\
1991 & 高维机制 & Bourgain 建立高维极大估计与限制猜想的联系\upcite{bourgain1991} \\
1995 & 发刷界 $(n+2)/2$ & Wolff 的“发刷”方法在 $\mathbb{R}^4$ 给出维数下界 3\upcite{wolff1995} \\
1999 & 加法组合入场 & Bourgain 把线关联与和集、差集估计挂钩\upcite{bourgain1999} \\
2000 & 三结构分类 & Katz--\L aba--Tao 识别粘性、平面性、颗粒性三种极端结构\upcite{klt2000} \\
2001 & 中维改进 & \L aba--Tao 在 $n\ge 4$ 把 Minkowski 下界推过 $(n+2)/2$\upcite{labatao2001} \\
2002 & 和差集迭代 & Katz--Tao 发展迭代式算术方法\upcite{katztai2002} \\
2006 & 多线性化 & Bennett--Carbery--Tao 建立多线性 Kakeya 定理\upcite{bct2006} \\
2009 & 多项式方法 & Dvir 解决有限域挂谷猜想\upcite{dvir2009} \\
2010 & 多线性端点 & Guth 给出端点形式\upcite{guth2010} \\
2015 & 多项式划分 & Guth--Katz 的多项式划分技术成形\upcite{guthkatz2015} \\
2018 & 多项式 Wolff 公理 & Guth--Zahl、Zahl、Katz--Rogers 联合给出 $\mathbb{R}^4$ 下界 $3+1/40$\upcite{guthzahl2018,zahl2018,katzrogers2018} \\
2021 & 平面刷 & Katz--Zahl 在 $\mathbb{R}^4$ 把 Hausdorff 下界推到 3.059\upcite{katz2021} \\
2022 & 粘性情形 & Wang--Zahl 证明粘性挂谷猜想\upcite{wangz2022} \\
2024 & Assouad 维数 & Wang--Zahl 给出 $\mathbb{R}^3$ 的 Assouad 维数定理\upcite{wangz2025assouad} \\
2025 & 三维全解 & Wang--Zahl 的 127 页论文解决三维集合猜想\upcite{wangz2025} \\
2026 & 三维简化 & Guth--Wang--Zahl 给出流线化证明\upcite{gwz2026} \\
2026 & OpenAI 手稿 & 三维极大函数与四维 Hausdorff 维数（后续章节讨论）\upcite{openai3d2026,openai4d2026} \\
\bottomrule
\end{tabularx}
\end{table}

时间线上有两个转折值得单独指出。第一个是 1971 年前后：Fefferman 的反例把“管族重叠”这个几何对象与球乘子这类分析算子的有界性绑在一起\upcite{fefferman1971}，从此挂谷问题不再只是组合几何的智力游戏，而成为调和分析的基础设施问题。第二个是 2009 年前后：Dvir 用多项式方法解决有限域版本\upcite{dvir2009}，随后 Guth--Katz 的多项式划分技术在欧氏空间“落地”\upcite{guthkatz2015}，到 2018 年演化为多项式 Wolff 公理体系\upcite{guthzahl2018,katzrogers2018}。这两次转折都不是挂谷内部自生的——第一次借的是 Fourier 分析的问题意识，第二次借的是有限域与组合几何的多项式方法。

\subsection{门派图谱}

\begin{table}[t]
\centering
\caption{挂谷问题相关门派与核心工具}
\label{tab:schools}
\footnotesize
\begin{tabularx}{\textwidth}{@{}L{2.6cm}L{4.4cm}Y@{}}
\toprule
\rowcolor{lightbg}\textbf{门派} & \textbf{核心工具} & \textbf{代表工作与角色} \\
\midrule
组合几何与几何测度 & 发刷、投影、Wolff 公理、近失（near-miss）反例 & Wolff\upcite{wolff1995}；Heisenberg 群、$SL_2$ 与二次超曲面反例界定了各方法的能力边界\upcite{zahl2025survey} \\
调和分析/极大算子 & 极大函数、限制估计、多线性 Kakeya & Fefferman、C\'ordoba、Bourgain、Bennett--Carbery--Tao；Bourgain--Guth 的多线性$\to$线性论证\upcite{fefferman1971,cordoba1977,bourgain1991,bct2006,bourgainguth2011} \\
加法组合 & 和差集迭代、Balog--Szemer\'edi 型定理 & Bourgain 1999 与 Katz--Tao 2002 把算术结构引入管族分析\upcite{bourgain1999,katztai2002} \\
多项式方法 & 多项式划分、Wolff 公理、Severi 型代数几何 & Dvir 2009 起点；Guth--Zahl、Zahl 把它做成 $\mathbb{R}^4$ 的主力工具\upcite{dvir2009,guthzahl2018,zahl2018} \\
极值与加权纲领 & 粘性、阴影（shading）、加权计数、极值反证 & Katz--Zahl 平面刷、Wang--Zahl 粘性与加权纲领、Guth--Wang--Zahl 流线化\upcite{katz2021,wangz2022,wangz2025,gwz2026} \\
分形几何与熵投影（邻近门派） & 熵平均、投影定理、Furstenberg 型估计、Fourier 衰减 & Hochman--Shmerkin、Ren--Wang、Orponen 系列；此前主要服务距离集与投影问题\upcite{hochmanshmerkin2012,renwang2025,orponend2024} \\
\bottomrule
\end{tabularx}
\end{table}

表 \ref{tab:schools} 的最后一栏是本文的一个铺垫：分形几何与熵投影这一门派，长期与挂谷问题相邻但并未深度介入主线。它们研究的是“集合在投影、距离、卷积之后还剩多少复杂度”，工具是熵与 Fourier 衰减；挂谷主线关心的是“方向的几何如何制造重叠”。两边的对象不同、语言不同，交流不多。对希望进入此题的读者，Wolff 的讲义至今仍是门内的标准入门读物\upcite{wolffbook2003}；OpenAI 手稿最引人注目的一点，正是把前一个门派的语言大规模翻译进挂谷问题——细节留到第 7 章。

\subsection{三维是怎么被拿下的}

三维的胜利来自极值纲领与结构分类的合流。Katz--\L aba--Tao 早就指出，逼近极端的管族要么在跨尺度上保持成组（粘性），要么在交点附近近似共面（平面性），要么并集组织成薄片（颗粒性）\upcite{klt2000}。Katz--Zahl 后来把三维 Hausdorff 下界推过 $5/2$\upcite{katz2019}，为粘性纲领补上了最后的中间台阶。此后二十年，业界始终无法把这三类结构做成一整套证明。

Wang--Zahl 分三步走完了这段路。2022 年先攻粘性情形\upcite{wangz2022}，2024 年拿下 Assouad 维数版本\upcite{wangz2025assouad}，2025 年 2 月的 127 页论文同时解决三维的 Hausdorff 与 Minkowski 维数\upcite{wangz2025}。随后 Guth--Wang--Zahl 给出流线化证明，把关键结构提炼为“凸聚类”条件下的并集估计——这条估计后来成为 OpenAI 两篇稿件共同的入口\upcite{gwz2026}。

三维的胜利有一个容易被忽略的特征：它的核心工具（凸聚类、棱柱、阴影、多尺度归纳）几乎全部来自门内。Nets Katz 评价 2025 年的结果“百年一遇”\upcite{guangming2026}，而这个结果本身是挂谷研究传统内部工具链成熟的产物。

\subsection{四维的现状：下界演进}

四维的情形完全不同。表 \ref{tab:r4bounds} 汇总了手稿出现之前的下界演进。

\begin{table}[t]
\centering
\caption{$\mathbb{R}^4$ 挂谷问题的下界演进（手稿发布前）}
\label{tab:r4bounds}
\footnotesize
\begin{tabularx}{\textwidth}{@{}L{2.4cm}L{2.2cm}L{2.2cm}Y@{}}
\toprule
\rowcolor{lightbg}\textbf{年份/工作} & \textbf{Hausdorff 下界} & \textbf{极大参数下界} & \textbf{方法} \\
\midrule
Wolff 1995\upcite{wolff1995} & $3$ & —— & 发刷 \\
\L aba--Tao 2001\upcite{labatao2001} & $>3$（上 Minkowski） & —— & 多尺度几何结构 \\
Guth--Zahl 等 2018\upcite{guthzahl2018,zahl2018,katzrogers2018} & $3+\frac{1}{40}\approx 3.025$ & —— & 多项式 Wolff 公理 \\
Katz--Zahl 2021\upcite{katz2021} & $3.059$ & $3.049$ & 平面刷 \\
Borges 等 2025\upcite{borges2025} & —— & $\frac{159+\sqrt{145}}{56}\approx 3.0543$ & 平面刷 + 两端 Furstenberg 估计 \\
Rai Choudhuri 2026（粘性集）\upcite{raichoudhuri2026} & $\frac{13}{4}=3.25$（限粘性） & —— & 粘性纲领 \\
OpenAI 手稿 2026\upcite{openai4d2026} & 声称 $4$ & ——（未触及四维极大） & 二次图卡 + 地平线 \\
\bottomrule
\end{tabularx}
\end{table}

两个细节值得留意。其一，四维 Hausdorff 下界从 3 爬到 3.059 用了二十六年，改进幅度不到 2\%，说明门内手段接近枯竭；其二，粘性情形的 3.25 是条件性结果，对任意挂谷集无效\upcite{wangz2025,raichoudhuri2026}。手稿若成立，等于把停在小数点后第二位的量直接推到理论上限，这一跨度的可信度评估必须结合验证状态一并看（第 9 章）。

\subsection{一条贯穿全史的现象}

回看表 \ref{tab:timeline}：加法组合来自 Bourgain 的算术视角\upcite{bourgain1999}，多项式方法来自有限域与组合几何\upcite{dvir2009}，多线性估计来自限制理论\upcite{bct2006}。挂谷问题每一次大踏步，背后都是一次跨门派借入。这条规律给第 7 章的谱系分析提供了参照系：OpenAI 手稿引入熵、耦合与实代数几何工具，并非异类，而是延续了同一个传统；问题只在于，这次借入的工具距离门内语言更远，消化成本更高。

\begin{keybox}[可记忆主句]
三维的解法是勾稽传统工具的集大成之作；四维的下界几十年只挪了小数点后两位，而防线之间的距离悬殊，正是 OpenAI 手稿引发关注的根源。
\end{keybox}
\section{OpenAI 工作的核验档案}

本章不谈数学，只做一件事：把与两篇稿件有关的事实固定下来，并把“未验证”写成可核对的条目，而不是一句笼统的保留意见。

\subsection{发布事实与时间线}

\begin{table}[t]
\centering
\caption{发布与修订时间线（2026 年 9 月--10 月）}
\label{tab:timeline-release}
\footnotesize
\begin{tabularx}{\textwidth}{@{}L{2.6cm}Y@{}}
\toprule
\rowcolor{lightbg}\textbf{时间} & \textbf{事件} \\
\midrule
2026-09-23 & 三维挂谷极大函数稿件成文（97 页，作者栏 OpenAI）\upcite{openai3d2026} \\
2026-09-24 & 四维挂谷集稿件成文（175 页）；另有一篇球面 Fourier 限制稿件同日成文，后被三维篇引为独立路径\upcite{openai4d2026,openaielliptic2026} \\
2026-10-06 & 官方发布“Sharing AI progress in mathematics”，仓库 openai/math 上线，共 372 个结果族；发布前曾咨询 IAS 的数学与人工智能咨询组（AGMAI）\upcite{openairelease2026,agmai2026} \\
2026-10-07 & 修订记录：撤回 3 篇（同一符号错误连带两篇依赖文献）、修订 14 篇、新增形式化 6 项；目录更新为 719 篇、Lean 覆盖 300/719\upcite{openairepo2026} \\
2026-10-07 & AHM（Association for Human Mathematics）发布声明并由陶哲轩博客转载，公开质疑发布方式\upcite{ahm2026,tao2026} \\
\bottomrule
\end{tabularx}
\end{table}

关于发布规模，仓库说明给出的是“719 篇手稿、372 个结果族，约 42\% 的顶层结果完成 Lean 形式化；平均每个结果消耗约 3 小时 ChatGPT Pro 思考算力；评测过程共投喂约 4000 道题”\upcite{openairepo2026}。这次放量发布并非首次：同年 5 月已有单位距离猜想否证的先例，8 月的十项进展则首次配套了 Lean 形式化\upcite{openaiunit2026,openaiten2026}。仓库同时释放了 10 份推理摘要，覆盖的是乘法函数两点关联、$\pi$ 的无理性指数等族，不包括第 074 族。仓库还披露两类例外：黎曼 $\zeta$ 函数零点自由区的写作经过人工润色，Hodge 猜想（CM 型阿贝尔簇）的结果不走统一流程\upcite{openairepo2026}。

\subsection{两篇稿件的参数与主定理}

\begin{table}[t]
\centering
\caption{第 074 族两篇稿件的基本参数（据稿件与目录信息整理）}
\label{tab:ms-params}
\footnotesize
\begin{tabularx}{\textwidth}{@{}L{2.5cm}L{2.6cm}Y@{}}
\toprule
\rowcolor{lightbg}\textbf{项目} & \textbf{三维篇} & \textbf{四维篇} \\
\midrule
标题 & The Kakeya maximal conjecture in three dimensions & Every four-dimensional Kakeya set has full Hausdorff dimension \\
成文日期 & 2026-09-23 & 2026-09-24 \\
篇幅 & 97 页 & 175 页 \\
主定理 & $K_\delta$ 把 $L^3(\mathbb{R}^3)$ 映入 $L^3(S^2)$，范数至多 $C_\varepsilon\delta^{-\varepsilon}$；等价地，阴影密度幂从 $\lambda^{K(\varepsilon)}$ 提升到 $\lambda^{3}$\upcite{openai3d2026} & 任意含每个方向单位线段的 $K\subset\mathbb{R}^4$ 满足 $\dim_{\mathrm{H}}K=4$，不设可测性、紧性、粘性假设\upcite{openai4d2026} \\
形式化状态 & 未列入 Lean 清单 & 未列入 Lean 清单 \\
评审状态 & 未经同行评审 & 未经同行评审 \\
\bottomrule
\end{tabularx}
\end{table}

三维篇的定位需要一句话说清楚：它解决的不是王虹--Zahl 已经拿下的集合维数问题，而是与之并列、更强的一档——极大函数猜想。用稿件自己的话说，Wang--Zahl 定理给出的阴影密度依赖是某个幂 $\lambda^{K(\varepsilon)}$，“换成 $\lambda^{3}$ 的需要”在他们定理之后就被明确写了出来\upcite{openai3d2026}。三维篇声称把这一步补齐；四维篇随后把三维篇的一条加权引理当作输入。两篇稿件在目录中被编为同一族，依赖关系见图 \ref{fig:deplink}（见附录 A）。

\subsection{依赖链：为什么两篇稿件必须一起评估}

\begin{evbox}[证据：依赖链的关键节点]
四维篇导言明确记载：归约中使用的三维估计是“OpenAI [24, Lemma 2.3] 的加权全时段板条界”，它在三维篇中由 Guth--Wang--Zahl 并集估计出发、经加权凸聚类证明得到；四维篇将其记为 Lemma 2.15。除此之外，四维篇还独立引入多元 Kakeya（行列式形式）、Wang--Zahl 的限制三元组点积定理、Orponen--de Saxc\'e--Shmerkin 的乘法卷积估计、图版本 Balog--Szemer\'edi--Gowers 论证与和集演算等输入\upcite{openai4d2026}。
\end{evbox}

这条链的后果直白：四维命题的成立，至少以三维篇的一条核心引理为前提；而三维篇又立在 Guth--Wang--Zahl 的并集估计之上\upcite{gwz2026}。核验四维稿件时，最省力的入口不是它的第 6 至 10 章（三区间分析），而是先把“三维加权板条引理”这一段单独读通——这也是本文附录 B 核验清单把该引理排在第一优先级的原因。

\subsection{验证状态：数字化读数}

\begin{evbox}[证据：可核对的三组读数]
（i）仓库的 Lean 形式化清单中不存在任何 Kakeya 相关条目，第 074 族两篇稿件均未形式化\upcite{openairepo2026}；（ii）同一批次已有撤回先例——3 篇 Hodge 相关稿件因一个符号错误导致稳定化迹抵消论证失效而撤回，其中两篇是因依赖该错误结果连带撤回，说明未形式化的稿件确实可能带着实质性缺陷通过内部流程\upcite{openairepo2026}；（iii）人类侧对照：三维挂谷猜想（Wang--Zahl）的证明已由中国科学院数学与系统科学研究院 MMAT 团队与开源组织 Project Numina 等合作完成 Lean 4 形式化\upcite{numina2026}——同一题材、同等量级的证明可以被形式化，这使 074 族的空白无法用“技术不可行”来解释。
\end{evbox}

三组读数指向同一个结论：074 族的验证负担完全落在人类数学家身上，且验证对象是一篇 175 页、自造术语密集、与另一篇 97 页稿件耦合的文本。稿件自身对这一风险的措辞是坦率的：“部分未形式化的结果可能存在问题，我们会尽快修复”\upcite{openairepo2026}。

\subsection{发布引发的争议（事实记录）}

AHM 的声明是目前最系统的一份反对文本，其要点为：数学家没有要求这类工作；在 AGMAI 的初始建议明确反对拿内部模型测试高级难题的前提下，OpenAI 仍以咨询组作为合法性来源；一次发布 700 余份文件“展示的不是学术，而是权力”；呼吁数学家停止与 OpenAI 合作、回到以人类理解为中心的科学愿景\upcite{ahm2026,tao2026}。声明由陶哲轩博客以客座帖转载，他本人既未署名、也未表态支持；帖子以下的讨论呈现明显的两极（有人赞赏结果、有人质疑发布方式），这场争论本身在第 9 章还会回到。

OpenAI 一侧的口径是：发布前咨询 AGMAI 并采纳其公开建议；为稿件建立修订与引用协议；为部分证明提供 Lean 形式化；将资助系列研讨会帮助学界理解这些结果\upcite{openairelease2026,agmai2026}。在其更早的“十项进展”发布中，公司对责任边界的表述是：数学论证由系统生成，他们负责整理与形式化，并为正确性负责\upcite{openaiten2026}。

\begin{cautbox}[边界：三件不能混淆的事]
第一，“未形式化”不等于“错误”，但它改变了举证责任的分布——主张成立的一方需要提供可核验的路径；第二，“未见公开错误”不等于“已验证”——175 页文本中，公开的专业审读记录目前几乎为零；第三，发布方的快速修订机制（撤稿与补丁）说明流程具备纠错能力，但纠错发生在发布之后，与“本文写作完成时结论是否成立”无关。
\end{cautbox}

\begin{keybox}[可记忆主句]
这份档案的核心只有一句话：两篇稿件、272 页、零形式化、零审评，且四维篇与三维篇共用一条未经核验的引理——所有关于四维挂谷的讨论，都必须从这里出发。
\end{keybox}
\section{四维证明的技术骨架}

四维篇的证明是一个反证：假设存在 Hausdorff 维数小于 $4$ 的挂谷集，构造一个带正“时间代价”的加权问题，再从这个问题中逼出一个自相矛盾的“改进”。本章按证明自身的模块顺序重建骨架：归约（第 2 章）、比较与延拓（第 3 章）、极值临界路径（第 4 章）、三个窄度区间（第 6、8--10 章）、收口与推论（第 10、11 章）。术语第一次出现时附原文，便于对照稿件。

\subsection{归约：把任意集合问题变成加权线族问题}

归约的定理陈述可以一字不改地引用\upcite{openai4d2026}：

\begin{theorem}[Reduction, 四维稿件 Theorem 2.17，转述]
若 $K\subset\mathbb{R}^4$ 不含任何可测性或紧性假设，包含每个方向的单位线段，且 $\dim_{\mathrm{H}}K<4$，则存在 $d<1$、$S,\eta>0$，以及一个属于容量类 $\mathcal{C}(d)$ 的温和齐次精确问题 $E$（版本 2），使得 $H(E)>0$。
\end{theorem}

这里连续出现了三个机制性概念，逐个解释。

\textbf{加权模型}。把“集合含每方向线段”翻译成概率语言：取一族轨迹（trajectory）$(t,\; y_0+tV,\; w_0+tW)$（其中 $y\in\mathbb{R}^2$ 是基础坐标，$w$ 是隐藏坐标），每条轨迹服从某个先验测度，时间独立均匀；再用一个入射选择（incidence selection）标记哪些“轨迹--时间”对被保留。这样处理后，问题从“集合的覆盖数”转成“加权入射测度的分布”，密度、角度集中度这类量都可以用指数记账。归约的第一步用 Hausdorff 覆盖的外测度可数次可加性，挑出一个斜率正测度、截距有界的子族——这一步不需要可测性，是“任意集合”能进入框架的关键。

\textbf{观察与密度}。在深度 $r$ 上做分箱观察：保留精确的基础坐标，砍掉隐藏坐标（记为 $[w]_{a_r/B}$），得到一组嵌套的观测 $p_r$。入射密度沿深度按 $n(r)=N_0^{-f(r)}$ 退化。反证假设 $\dim_{\mathrm{H}}K<4$ 体现在密度退化速度上：存在 $d<1$，使 $f(L)-f(r)\le d(L-r)$，即深度每推进一格，密度损失不到一个完整幂（容量类 $\mathcal{C}(d)$ 的第一条）。第二条是角度非集中：隐藏观察细化时，速度方向的分布不许过度集中。

\textbf{图卡、覆盖与地平线}。图卡（chart）是证明的局部对象：一个移动的空间盒、一个时间区间 $I$，加上一个二次多项式测试 $P$，要求 $P$ 在足够多的被选轨迹上保持小（厚度受控）。图卡自身带质量要求，用“选定的时间轴对”度量，防止“只拟合了可忽略的例外族”被当成进展。区间长度写成 $N_0^{-h+o(1)}$ 时，$h$ 叫时间代价（time cost）；一个图卡系统的覆盖量是其选定入射质量，质量 $N_0^{-o(1)}$ 叫实质覆盖。于是：

\begin{equation}
H(E)\;=\;\inf\{\, h \;:\; \text{存在带实质覆盖的真终端图卡系统，区间长度} = N_0^{-h+o(1)} \,\},
\end{equation}

即地平线（horizon）——在“拟合到最终厚度”这件事上，时间局部化至少要付的代价。归约的结论是 $H(E)>0$：反例逼出了一笔正代价。（地平线的形式定义、图卡与真图卡的完整条件、以及归约定理的形态，转录于附录 C.1，供逐条核对。）

\begin{keybox}[归约的要点]
反证假设（维数亏空）被换算成一条关于指数的正下界（地平线为正）。此后全部工作变成：证明 $H(E)=0$，或者证明 $H(E)>0$ 会导致矛盾。整个四维证明的剩余部分都在执行后一种路线。
\end{keybox}

\subsection{比较原理：跨尺度把局部描述接起来}

有了图卡，下一步是把不同尺度上的局部描述拼接起来。第 3 章的比较原理（comparison of algebraic sheets）处理的正是拼接的合法性：多项式测试的局部图像给出“标量叶片”（scalar sheets），即多项式方程的解析分支（精确定义与 Lemma 3.2、Lemma 3.3 的定量陈述见附录 C.1）。密度约束与保留下来的低界质量给出一个关键定量事实：在精度 $\delta$ 上，每个局部描述只需一个大小至多 $\delta^{-d+o(1)}$ 的列表——由于 $d<1$，列表规模小于任何完整幂。

比较原理的机制值得复述：两个这样的列表不可能“频繁地取值一致、同时保持换算后的导数数据分离”；否则，交替抽取新条件下的叶片与时间，会把持续的导数失配拼成一个多项式端点映射，其像的体积超过密度约束允许的上限。成功路径保持为非归一化子测度，因此体积论证保留了对先前选择的依赖——这个技术细节保证比较不会悄悄放宽假设。

\begin{evbox}[证据：d<1 与比较原理的耦合]
容量类 $\mathcal{C}(d)$ 要求 $d<1$，稿件给出的原因正是比较原理：列表规模 $\delta^{-d+o(1)}$ 只有在 $d<1$ 时才能被“体积超界”论证压住\upcite{openai4d2026}。这类“一个常数被证明机制的结构所选定”的安排，是手稿设计感的典型体现。
\end{evbox}

\subsection{极值选择与临界路径}

第 4 章把问题摆到一个近乎极值的加权问题上：沿着嵌套图卡行进，相对密度与时间剖面会线性化，出现正整数的时间速率 $k>0$；再做一次极值选择，控制“较小的空间轴能比时间尺度细多少”，把最优窄度（optimized narrowness）记为 $\ell$。$\ell$ 大体上衡量管族压扁的方向——$\ell=0$ 是各向同性，$\ell$ 有限是有限窄度，$\ell=\infty$ 是无界窄度。

“改进”（improvement）的定义由此而来：要么在保持覆盖的情况下把时间代价严格推低，要么在 $\ell$ 有限时于固定时间上换取额外的窄度增益。极值选择保证：任何一处这样的改进，回到原问题后都违反极值性。于是证明的战术目标变成了“在每一个实质性残余上制造一次改进”。

从局部候选到“实际改进”的转换是第 4 章的另一半工作，收在命题 Proposition 4.11 与引理 Lemma 4.12：先用有限贪心赋值把候选情景压实成有实质覆盖的单一情景，再用小列表比较把候选测试组装成改进的图卡测试；最后 Lemma 4.12 把有限列表从分析用子测度回填到原测度（误差由基网格上的全变差控制）。这一步是整套证明里最容易“看起来像行政细节、实际是逻辑要冲”的地方——没有它，局部拟合无法升级为对原问题成立的改进。

\subsection{三个窄度区间}

\begin{table}[t]
\centering
\caption{三个窄度区间：假设、机制与产出}
\label{tab:regimes}
\footnotesize
\begin{tabularx}{\textwidth}{@{}L{1.9cm}L{3.4cm}Y@{}}
\toprule
\rowcolor{lightbg}\textbf{区间} & \textbf{假设} & \textbf{机制与产出}（依据稿件第 6、8--10 章） \\
\midrule
$\ell=0$ 各向同性 & 分离速度足够多 & 沿单线的条件标量估计给出多项式拟合；多元 Kakeya 与插值合并；换线迫使近似守恒，产出势函数；宽带方向的圆锥障碍由隐藏坐标编码三次、四次项修回二次测试（次数恢复，Proposition 6.8） \\
$0<\ell<\infty$ 有限窄度 & 投影后部分区间足够宽 & 宽投影区间直接由比较原理改进；否则平面刚性把轨迹压入水平模型 $Y'=-Z+tv$，自然盒型为水平尺度 $H$、横向 $H^{2}$；临界速率 $2k=1$ 时三条刷新线段填满三维基础体积，偏离临界速率时回归宽区间，$2k<1$ 时先做自举 \\
$\ell=\infty$ 无界窄度 & 窄度趋于无穷 & 两个嵌套图卡标签把轨迹在全归一化时间区间上局部化；归一化造成的密度指数大漂移在相对比值中抵消；最终标量投影产出时间代价趋于零的图卡，与正地平线矛盾 \\
\bottomrule
\end{tabularx}
\end{table}

三个区间共享一条供给线：各向同性情形的标量版本（单空间坐标）先被证明“零地平线”，其副产品是标量投影定理（Theorem 8.2），供另外两个区间做降维使用\upcite{openai4d2026}。这种“先在最难区间把一维辅助模型做出来，再让其他区间站上去”的安排，使三区间不是三条平行线，而是一棵树。

表 \ref{tab:regimes} 中两处机制值得展开一句。其一，圆锥障碍与次数恢复：各向同性情形产出的势函数初始次数可能高于 $2$，而“宽带方向落在圆锥附近”是唯一无法直接排除的障碍形态；稿件识别出残余的三次、四次项，用原本闲置的隐藏坐标把它们编码进二次修正，从而恢复图卡所需的次数。其二，水平模型与临界速率：$Y'=-Z+tv$ 这一形式的水平盒为 $H\times H^{2}$，其“非交换”的水平方向迫使标量叶片近似仿射，从而让 $2k=1$ 的三腿切换恰好把三维基础体积填满。

\subsection{收口、几何推论与三维范式的对照}

合成段（第 10 章）把三个区间的产出拼回原问题：每个区间先在加权族的实质残余上产出局部候选，再经贪心覆盖与小列表比较升级为可回填到原问题的改进图卡系统，与极值性矛盾。归纳法与 $H(E)>0$ 的阳性下界共同给出 $\dim_{\mathrm{H}}K\ge 4$；上界是空间的平凡上界，定理闭合\upcite{openai4d2026}。

第 11 章把主定理摊成一组推论：有界挂谷集同时具有满的下、上 Minkowski 维数与 packing 维数；对 $d\ge 4$，经四维子空间的正交投影可得 $\dim_{\mathrm{H}}E\ge 4$（$d=5$ 也包含在内，但离 $5$ 的完整猜想还差得远）；任意方向集版本 $\dim_{\mathrm{H}}E\ge 1+\dim_{\mathrm{H}}D$ 经 Keleti--M\'ath\'e 的转移得到；常曲率流形上的 Nikodym 集与固定平移不变相下的弯曲挂谷族，经 Gao--Liu--Xi 的转移定理得到。这些推论都附带各自的假设条件，稿件据此逐条列出\upcite{openai4d2026,keletimathe2023,gaoliuxi2025}。

\begin{table}[t]
\centering
\caption{四维稿件与三维范式在对象层面的对照}
\label{tab:3d4d}
\footnotesize
\begin{tabularx}{\textwidth}{@{}L{2.8cm}L{4.6cm}Y@{}}
\toprule
\rowcolor{lightbg}\textbf{维度} & \textbf{结构分类} & \textbf{证明载体} \\
\midrule
三维（Wang--Zahl 传统） & 粘性、平面性、颗粒性（近极端管族的三种形态） & 棱柱与阴影加权、多尺度归纳、凸聚类；关键定理在四维失效（二次曲面可容纳三维线族） \\
四维（OpenAI 手稿） & 按窄度参数 $\ell$ 划三区间，取代形态分类 & 二次图卡与质量约束、地平线极值反证、标量投影降维、水平模型 \\
\bottomrule
\end{tabularx}
\end{table}

表 \ref{tab:3d4d} 是本文对两代方法差异的概括（本文判断，\cfM）：三维的范式是“先分类极端形态，再逐一击破”，四维稿件的范式是“不分类形态，改为对局部拟合的时间代价做极值分析”——后者对窄度参数 $\ell$ 的连续性处理，正是为了绕开四维中比三维复杂得多、可能穷举不尽的聚集形态。

\begin{keybox}[可记忆主句]
四维证明的骨架一句话概括：反例逼出一条正地平线，三区间各自证明能在实质残余上制造一次改进，而改进与极值性不相容——于是反例不能存在。
\end{keybox}

\subsection{三维姊妹篇的骨架（作为对照）}

三维篇的目标是把阴影密度幂补到 $\lambda^{3}$。其机制与四维篇同源但更紧凑：先建立一个对密度敏感的临界指数 $h(p,z)$，把“极大定理成立”等价为“临界指数为零”；再假设临界指数为正，构造取等序列（equality configurations），用局部质量界与多元 Kakeya 逼出“短块上方向近共面”，用平面 Furstenberg 估计与控制支撑代价的板条 Nikodym 估计把“时间亏空”压成零；随后证明取等配置在跨尺度上必然收敛到一条规范时间剖面 $F(s)=\beta(s-\tau)_{+}$，并在同一指标的两个事件（双帧实验）之间建立投影约束，与熵需求矛盾\upcite{openai3d2026}。这条链条中，“同一根管子取两个时刻”的双帧实验、以及对钉住（pinning）估计的改造，与熵账本一起构成了第 7 章谱系分析的主要对象。

\begin{keybox}[可记忆主句]
三维篇与四维篇共享同一套记账语言（时间代价、覆盖、极值改进），差别在于三维篇把矛盾压在“双帧投影约束”上，四维篇把矛盾摊在“三区间改进”上。
\end{keybox}
\section{技术谱系（一）：直系继承与门内改造}

第 5 章重建了“怎么证”，本章回答“从哪来”。用户关心的问题是：在手稿使用的全部核心技术里，哪些是挂谷研究门派内部已有的资产，哪些是门派没有用过的东西；哪些跨了门派、跨了子学科。先把分类法与证据规则立好，再逐条归档。

\begin{keybox}[四分类法]
\textbf{A 直系继承}：门内成熟工具，角色基本不变或仅有加强；\textbf{B 门内改造}：能找到门内前身，但对象或角色被结构性替换；\textbf{C 跨门派借入}：工具来自挂谷传统之外，需要翻译才能使用；\textbf{D 谱系内首次系统使用}：以现有检索范围未见等价先例。A、B 归本章，C、D 归第 7 章。
\end{keybox}

证据规则有三条，全部可复核。第一，以手稿自述为准绳：稿件在标明来源时使用 adapts、built on、uses 之类的措辞，本文直接采信并标注；第二，以门内综述（Zahl 2025 挂谷综述）描述的“工具箱清单”为基线，判断某工具是否属于门派内标准装备\upcite{zahl2025survey}；第三，对门内两篇主力文献（Wang--Zahl 127 页、Guth--Wang--Zahl 47 页）做全文关键词扫描，用词频对照为“未使用”类断言提供量化证据（扫描结果见第 7.3 节）。三条规则解决不了的情形（例如某工具以别名出现、或只在小范围讨论中出现过），一律降低置信度并写明理由。

\subsection{直系继承清单}

\begin{table}[t]
\centering
\caption{直系继承的门内资产}
\label{tab:inherit}
\footnotesize
\begin{tabularx}{\textwidth}{@{}L{3.1cm}L{4.2cm}Y@{}}
\toprule
\rowcolor{lightbg}\textbf{门内资产} & \textbf{来源} & \textbf{在手稿中的角色} \\
\midrule
发刷方法与近失反例传统 & Wolff 1995；Heisenberg、$SL_2$、二次超曲面反例\upcite{wolff1995,zahl2025survey} & 划定四维障碍的形状（二次曲面可载三维线族），为“二次图卡”提供问题意识 \\
多项式 Wolff 公理谱系 & Guth--Zahl、Zahl、Katz--Rogers\upcite{guthzahl2018,zahl2018,katzrogers2018} & 直接继承为“对低次代数曲面聚集的控制”这一主题；对象层面经改造（见 6.2 节） \\
Guth--Wang--Zahl 并集估计 & 流线化证明的第 1.1 节估计\upcite{gwz2026} & 经三维篇的加权板条引理进入四维归约，是依赖链的根 \\
多元 Kakeya（行列式形式） & Bennett--Carbery--Tao、Guth、Carbery--Valdimarsson\upcite{bct2006,guth2010,carberyvaldimarsson2013} & 各向同性区间的方向独立性控制与组合工具 \\
平面刚性输入 & Wang--Zahl 限制三元组点积定理；Ren--Wang 平面 Furstenberg 估计（Wang--Wu 形式）；Orponen--Shmerkin--Wang 投影估计\upcite{wangz2022,renwang2025,wangwu2024,orponensw2024} & 三维篇的平面填充与四维篇的平面刚性（“保边”改造见第 7 章） \\
阴影与加权计数语言 & Katz--Zahl 平面刷的 shadow 计数；Wang--Zahl 的 shading 体系\upcite{katz2021,wangz2022} & 三维篇仍用阴影语言（全文 48 次），四维篇转写为“入射选择 + 覆盖量”的测度语言 \\
棱柱—板条工具箱 & Wang--Zahl 的 prism 体系（全文 335 次）→ Guth--Wang--Zahl 的 plank（71 次）\upcite{wangz2025,gwz2026} & 三维篇的“全时段板条 + 加权聚类”（82 次），再传入四维篇 \\
\bottomrule
\end{tabularx}
\end{table}

表 \ref{tab:inherit} 里藏着一个值得单独指出的传递现象：门内的空间记账单位经历了一条清晰的演化——\textbf{棱柱（prism）}在 Wang--Zahl 处是主力对象，到 Guth--Wang--Zahl 被\textbf{板条（plank）}取代（词频从 335 次降到 71 次），到 OpenAI 三维篇进一步变成\textbf{加权全时段板条}（plank 82 次、prism 归零）。四维篇拿到的，其实不是王虹--Zahl 的原始工具，而是经过 Guth--Wang--Zahl 流线化、再由 OpenAI 三维篇加权化之后的第三代资产。这份手稿在数学继承上的“血缘”相当近。

\subsection{门内改造清单}

\begin{table}[t]
\centering
\caption{门内改造：可辨认的前身与结构性替换}
\label{tab:modify}
\footnotesize
\begin{tabularx}{\textwidth}{@{}L{3.1cm}L{4.4cm}Y@{}}
\toprule
\rowcolor{lightbg}\textbf{改造项} & \textbf{门内前身} & \textbf{替换内容} \\
\midrule
三窄度区间 & Katz--\L aba--Tao 的粘性/平面性/颗粒性三形态分类\upcite{klt2000} & 形态穷举换为参数分层：按最优窄度 $\ell$ 划三区间，用极值选择导出一致的改进目标，规避“四维聚集形态可能穷举不尽”的风险 \\
二次图卡体系 & 多项式划分与多项式 Wolff 公理\upcite{guthkatz2015,guthzahl2018} & 全局“划分 + 排除坏集”换为局部“移动盒 + 二次测试 + 质量约束 + 可嵌套归一化”；对二次结构的利用从抑制改为拟合 \\
密度幂接力 & Wang--Zahl 定理后遗留的 $\lambda^{K(\varepsilon)}\to\lambda^{3}$ 问题\upcite{openai3d2026,wangz2025} & 三维篇按文献线索原样接管该问题，把“补一个尖锐幂次”当作独立目标完成 \\
极值机制的转写 & 近失反例的构造传统；粘性纲领的极值配置分析\upcite{zahl2025survey,wangz2022} & 把“构造反例测试方法极限”转写为“反例逼出正时间代价（地平线），再用改进与极值性对撞”的证明资源 \\
连续化记账 & 门内的离散尺度层级（dyadic 尺度上逐层归纳） & 时间代价 $h$ 与密度剖面 $f(r)$ 的连续指数记账（含 Lipschitz 剖面与容量类 $\mathcal{C}(d)$） \\
\bottomrule
\end{tabularx}
\end{table}

表 \ref{tab:modify} 中前两行是本章最重要的判断（\cfM）：四维篇对门内范式的替换，集中在“分类对象”与“拟合载体”两处。三维的粘性/平面性/颗粒性分类，本质上是对近极端管族做形态学穷举；这套方法在三维有效，是因为低维空间中可利用的代数退化形态有限。四维里能承载大量线族的代数集合骤然增多（二次超曲面只是开端），穷举路线前景暗淡，于是手稿改用参数分层：不问“形态是什么”，改问“最优窄度参数落在哪一段”。类似地，多项式划分把“低次多项式零点集”当作抑制聚集的工具（划分单元、排除坏集），而图卡把“二次多项式测试”当作拟合聚集的载体（在移动盒上让测试保持在允许厚度内）——同属多项式语言，姿态从防守转为进攻。

第三行的“接力”值得补充一个细节：三维篇导言明确写着，$\lambda^{K(\varepsilon)}$ 换成 $\lambda^{3}$ 的需要在 Wang--Zahl 定理之后被明确提出，并给出定位\upcite{openai3d2026,wangz2025}。这不是事后附会——遗留问题写在文献里，手稿把它捡起来做完了。这类“拾遗式”选题与四维篇的“跨越式”选题（从 3.059 直接声称到 4）性质不同，验证负担也不同：三维篇若成立，是一次完成度很高的接力；四维篇若成立，则是一次范式级的推进。

\subsection{继承的边界}

\begin{cautbox}[未继承的东西]
手稿把门内的数学资产继承得相当完整，却没有继承门内的验证标准。三维挂谷证明（王虹--Zahl）的 Lean 4 形式化已由中科院 MMAT 团队与 Project Numina 等合作完成\upcite{numina2026}，而 074 族两篇稿件在仓库形式化清单中不存在任何条目\upcite{openairepo2026}。“继承数学、不继承验证”这一点，是本文对谱系分析给出的第一个带判断的结论（\cfH）——它的后果在第 9 章展开。
\end{cautbox}

\begin{keybox}[可记忆主句]
这份手稿在门内的血缘比表面看起来更近：它继承的不是二十年前的发刷，而是刚刚流线化、又被自己加权化的第三代板条工具；真正的替换只发生在两处——分类对象与拟合载体。
\end{keybox}
\section{技术谱系（二）：跨门派借入与谱系内首次使用}

本章处理四分类法里最有信息量的两类：C 跨门派借入与 D 谱系内首次系统使用。前者回答“哪些技术是挂谷门派没有用过的”，后者回答“哪些构造成套出现在这份手稿里像第一次”。两条线索最后合并成 7.3 节的词频证据表——一张可以逐格复核的对照表。

\subsection{跨门派借入清单}

\begin{table}[t]
\centering
\caption{跨门派借入：来源、门内先前状态与手稿中的用途}
\label{tab:cross}
\footnotesize
\begin{tabularx}{\textwidth}{@{}L{2.9cm}L{2.6cm}L{2.6cm}Y@{}}
\toprule
\rowcolor{lightbg}\textbf{技术} & \textbf{来源门派/子学科} & \textbf{门内先前状态} & \textbf{在手稿中的用途} \\
\midrule
条件熵与互信息账本 & 信息论（Cover--Thomas）；分形几何的熵方法（Hochman--Shmerkin 谱系）\upcite{coverthomas2006,hochmanshmerkin2012} & 未见（门内两篇主力文献零命中） & 三维篇 §6--§7 的动帧熵增量、坐标熵率下界；四维篇 §5 高阶喷流拟合中的熵估计与密度账本 \\
概率耦合与乘积律比较 & 概率论 & 未见 & 三维篇 §8 双帧实验：实际联合律与乘积律作比较，约束失败则乘积律以多项式小概率逃出支撑，逼出正的互信息需求，与其上界矛盾 \\
投影钉住（pinning）bootstrap & 分形投影理论（Shmerkin--Wang；Orponen--Shmerkin--Wang）\upcite{shmerkinwang2025,orponensw2024} & 未见；手稿自述 adapts & 三维篇 §9 乘积斜率 $\Delta t\,\Delta\vartheta$ 的有限尺度估计 \\
乘法卷积的 Fourier 衰减估计 & 分形几何/Fourier 衰减（Orponen--de Saxc\'e--Shmerkin）\upcite{orponend2024} & 未见 & 四维篇 Proposition 7.6 之前的构造：标量系数的小区间概率，服务平面刚性 \\
图版本 Balog--Szemer\'edi--Gowers 与和集演算 & 极值组合/加法组合（Sudakov--Szemer\'edi--Vu；Tao--Vu）\upcite{ssv2005,taovu2006} & 非图版 BSG 已经 Bourgain 1999 内化；图版本及其“保留原始关联边”的用法未见 & 四维篇 §7 保边平面刚性：把拟合出的平面结构回填到原始轨迹的关联边上 \\
实代数几何的定量纪律 & 实代数几何/逼近论（Basu--Pollack--Roy 的半代数算法；Roytwarf--Yomdin 的 Bernstein 类）\upcite{bpr2006,roytwaryomdin1997} & 未见（门内多项式方法只用到次数界与 Bézout 型事实） & “温和数据”（tempered data）装置；端点比值的有界性论证；次数恢复中的多项式 Remez 迭代 \\
Nikodym/弯曲挂谷转移 & 几何测度与投影（Keleti--M\'ath\'e；Gao--Liu--Xi；Nadjimzadah）\upcite{keletimathe2023,gaoliuxi2025,nadjimzadah2026,nadjimzadah2026b} & 近年内新兴 & 两篇稿件的推论章：常曲率流形上的 Nikodym 集、固定相位下的弯曲挂谷族 \\
\bottomrule
\end{tabularx}
\end{table}

表 \ref{tab:cross} 的阅读方式建议从第三列看起：这一列是本文试图回答的“门派有没有用过”。七个条目里，四个此前几乎与挂谷主线无交集（熵账本、耦合、Fourier 衰减、实代数几何纪律），两个属于“同属组合数学但具体形态未见”（图版本 BSG、钉住改造），一个属于近年新工具的正常借用（转移定理）。

其中两处借用值得多写两句。其一，熵账本的性质与以往门内工具都不同：发刷、划分、板条都是关于“几何对象如何分布”的陈述，熵是关于“信息在坐标间如何携带”的陈述。三维篇把一根管子的两个时刻视为同一指标的两次观测，用条件熵与互信息预测“两帧之间的差异必须携带多少信息”，再把它与“可得信息量的上界”对撞。这是把物理学的“信息守恒”直觉形式化为不等式，在挂谷文献里此前没有这种笔法。其二，实代数几何的引入解决的是一类此前被回避的“脏活”：证明中反复需要“某个比值是一维半代数函数，因而有多项式增长界”这样的定量事实，以维持全局约定的温和性（多项式界的复杂度不爆炸）。门内传统对此通常选择绕开——用更粗的估计掩盖；手稿选择正面建立纪律，这是工程口味很重的选择（本文判断，\cfM）。

\subsection{谱系内首次系统使用的构造}

\begin{table}[t]
\centering
\caption{谱系内首次系统使用的构造成套（D 类）}
\label{tab:novel}
\footnotesize
\begin{tabularx}{\textwidth}{@{}L{3.0cm}XL{1.4cm}@{}}
\toprule
\rowcolor{lightbg}\textbf{构造} & \textbf{内容与检索依据} & \textbf{置信度} \\
\midrule
地平线框架 & 把“局部拟合精度”量化为时间代价并取极值资源（$H(E)>0$ 的逼出与“改进即矛盾”的对撞）。门内两篇主力文献与 Zahl 综述均无对应机制 & 高 \\
二次图卡体系 & 移动盒 + 时间区间 + 二次测试 + 质量约束（按选定时间轴对度量）+ 可嵌套归一化。词频佐证：chart 在门内文献 0 次，在两篇稿件分别 123、179 次 & 高（作为词汇与对象）／中（作为机制，与划分、Remez 存在亲缘） \\
代数叶片比较原理 & 用标量叶片的列表匹配与端点映射体积论证做跨尺度延拓；两色匹配控制失配。未见等价先例，警惕与“列表译码”类思想的潜在相似性 & 中 \\
次数恢复（圆锥障碍） & 识别势函数残余的三次、四次项，用闲置隐藏坐标编码为二次修正，恢复图卡所需的次数（Proposition 6.8） & 中高 \\
三区间与临界速率装置 & 水平模型 $Y'=-Z+tv$ 的 $H\times H^{2}$ 盒型；$2k=1$ 的三腿切换；$2k<1$ 的自举；无界窄度下密度指数漂移在比值中抵消 & 中高 \\
双帧投影约束实验（三维篇） & 同一标记载体取两个事件构造剪切变换，比较实际联合律与乘积律，配合乘积斜率钉住导出投影约束 & 中高 \\
保边回填（Lemma 4.12 类） & 把分析子测度上的有限列表回填到原测度，用基网格全变差控制误差；保证局部拟合升级为原问题上的改进 & 中 \\
\bottomrule
\end{tabularx}
\end{table}

表 \ref{tab:novel} 的置信度差异本身就是一份诚实的边界声明：词汇与对象层面的新颖（地平线、图卡）证据最硬——门内文献完全不使用这些词；机制层面的新颖（比较原理、次数恢复）只能靠机制结构与检索范围来支持，无法排除“某些技巧在未广泛传播的工作中出现过”。凡此含糊处，本文统一取较低一档置信度。地平线、图卡、两色匹配与改进候选的精确定义及命题形态，转录于附录 C.1，供逐条核对。

\subsection{词频证据表：门派工具箱的定量对照}

对四篇论文全文做小写化词频统计（Wang--Zahl 127 页、Guth--Wang--Zahl 47 页、OpenAI 三维篇 97 页、四维篇 175 页），结果如表 \ref{tab:wordfreq}。表头缩写：W--Z 为 Wang--Zahl 2025；G--W--Z 为 Guth--Wang--Zahl 2026；OA 三维、OA 四维为 OpenAI 两篇稿件。

\begin{table}[t]
\centering
\caption{关键技术在四篇文献中的词频对照（全文、小写化、大小写不敏感）}
\label{tab:wordfreq}
\footnotesize
\begin{tabularx}{\textwidth}{@{}Xrrrr@{}}
\toprule
\rowcolor{lightbg}\textbf{关键词} & \textbf{W--Z 2025} & \textbf{G--W--Z 2026} & \textbf{OA 三维} & \textbf{OA 四维} \\
\midrule
entropy（熵） & 0 & 0 & 65 & 33 \\
mutual information（互信息） & 0 & 0 & 4 & 3 \\
coupling（耦合） & 0 & 0 & 10 & 0 \\
pinning（钉住） & 0 & 0 & 15 & 0 \\
Bernstein（类） & 0 & 0 & 0 & 2 \\
semialgebraic（半代数） & 0 & 0 & 0 & 21 \\
quantifier elimination（量词消去） & 0 & 0 & 0 & 4 \\
Remez（多项式迭代） & 0 & 0 & 0 & 17 \\
Balog--Szemer\'edi（BSG） & 0 & 0 & 0 & 3 \\
Fourier decay（Fourier 衰减） & 0 & 0 & 0 & 1 \\
sumset（和集） & 0 & 0 & 0 & 7 \\
\midrule
chart（图卡） & 0 & 0 & 123 & 179 \\
sheet（叶片） & 0 & 0 & 0 & 33 \\
shading（阴影） & 144 & 77 & 48 & 1 \\
prism（棱柱） & 335 & 6 & 0 & 0 \\
plank（板条） & 0 & 71 & 82 & 24 \\
convex（凸） & 192 & 74 & 24 & 7 \\
\bottomrule
\end{tabularx}
\end{table}

表 \ref{tab:wordfreq} 给出三条读数。第一条：全部“外门派”词汇（entropy、mutual information、coupling、pinning、Bernstein、semialgebraic、quantifier elimination、Remez、BSG、Fourier decay、sumset）在门内两篇主力文献中命中为零，在 OpenAI 两篇中集中出现——这是“门派没有用过这些技术”最直接的语言层证据。第二条：门内主力词汇的迁移轨迹清晰可见——棱柱从 335 次衰减到 0，板条在流线化证明中成为新单位（71 次），阴影语言在两篇 OpenAI 稿件中逐渐退场（144$\to$77$\to$48$\to$1）。第三条：chart 与 sheet 两个词在门内为零、在稿件中高密度出现，说明图卡与叶片不是翻译过来的旧对象，而是稿件自造的工作词汇。

\begin{cautbox}[词频证据的限度]
词频测量的是语言，不是思想。它支持“某词汇体系未被门内使用”这类断言，却无法单独证明“某思想未被门内使用”；概念可能以别名出现在未进入本文检索范围的工作中。正因如此，表 \ref{tab:wordfreq} 只被用来支撑 C 类判断（表 \ref{tab:cross} 第三列），D 类判断（表 \ref{tab:novel}）另行采用机制结构对照，并保留更低一档置信度。
\end{cautbox}

\subsection{小结：一次距离更远的借入}

把本章与前章合起来看：手稿在门内继承的是刚出厂的三代资产（板条、流线化、加权），向外借的是信息论、分形几何与实代数几何三家的语言。门内历史上每次借入（加法组合、多项式方法、多线性估计）都花了一代人的时间才完成消化；本次借入的距离更远——从“几何对象的分布”跨到了“信息的携带”，从“次数界”跨到了“半代数复杂度纪律”。这既解释了手稿的体量（175 页），也预示了它的消化成本（第 9、10 章）。

\begin{keybox}[可记忆主句]
门派内部从未见过熵、耦合与半代数纪律进入挂谷问题的主线：这一次借入的不仅是工具，还有一整套记账方式与证明纪律。
\end{keybox}
\section{统一比较矩阵}

第 6、7 章分开讲了“继承什么”与“借入什么”。本章把这条研究线上的主要方法放进一张矩阵横向比较：发刷、多项式 Wolff 公理、平面刷、粘性纲领、三维全解、OpenAI 三维篇与四维篇。比较的四个轴是目标与维度、结论强度、假设、核心工具，另附跨门派成分与验证状态两列。

\begin{table}[t]
\centering
\caption{挂谷研究主要方法的统一比较矩阵}
\label{tab:matrix}
\scriptsize
\begin{tabularx}{\textwidth}{@{}L{2.0cm}L{2.3cm}L{2.5cm}XL{1.5cm}L{1.7cm}@{}}
\toprule
\rowcolor{lightbg}\textbf{方法（年份）} & \textbf{目标} & \textbf{结论强度} & \textbf{核心工具} & \textbf{跨门派成分} & \textbf{验证状态} \\
\midrule
发刷（1995）\upcite{wolff1995} & $\mathbb{R}^n$ 维数下界 & $(n+2)/2$；$\mathbb{R}^4$ 得 $3$ & 发刷、管交体积估计 & 低 & 已证、同行评审 \\
多项式 Wolff 公理（2018）\upcite{guthzahl2018,zahl2018,katzrogers2018} & $\mathbb{R}^4$ 维数下界 & $3+\frac{1}{40}$ & 多项式划分、Wolff 公理、Severi 型代数几何 & 低（有限域方法早已内化） & 已证、同行评审 \\
平面刷（2021）\upcite{katz2021} & $\mathbb{R}^4$ 维数与极大参数 & $3.059$ / $3.049$ & 平面刷、阴影计数、多尺度组织 & 低 & 已证、同行评审 \\
粘性纲领（2022--2026）\upcite{wangz2022,raichoudhuri2026,wangzakharov2026} & 粘性集的维数 & $\mathbb{R}^3$ 全维（限粘性）；$\mathbb{R}^4$ 粘性 $13/4$ & 多尺度归纳、加权计数、稠密阴影条件 & 低 & 已证（条件性结论） \\
三维全解（2025--2026）\upcite{wangz2025,gwz2026} & $\mathbb{R}^3$ 集合 & $=3$（Hausdorff 与 Minkowski） & 凸聚类、棱柱$\to$板条、阴影、加权并集估计 & 低 & 已证、同行评审（流线化版本 2026） \\
OpenAI 三维篇（2026）\upcite{openai3d2026} & $\mathbb{R}^3$ 极大函数 & $\lambda^{3}$（临界密度幂）；$L^3\to L^3$ 带 $\delta^{-\varepsilon}$ & 加权全时段板条 + 熵账本 + 耦合 + 钉住 + Furstenberg & 中高 & 未形式化、未评审（声称） \\
OpenAI 四维篇（2026）\upcite{openai4d2026} & $\mathbb{R}^4$ Hausdorff 维数 & $=4$（任意集合） & 二次图卡 + 地平线 + 标量投影 + 借入工具集 & 高 & 未形式化、未评审（声称） \\
\bottomrule
\end{tabularx}
\end{table}

\subsection{三条读表线索}

\textbf{线索一：假设在收缩。}从粘性纲领到三维全解，条件从“粘性、稠密阴影”收缩到无条件；四维篇承接同样姿态，明确不设可测性、紧性与粘性假设\upcite{openai4d2026}。假设收缩意味着结论覆盖面变宽，也意味着证明必须处理更野的配置——这是四维篇体量膨胀的直接原因。

\textbf{线索二：结论在精确化。}三维篇把遗留的密度幂 $\lambda^{K(\varepsilon)}$ 精化为 $\lambda^{3}$，属于“在已知框架内把最优点补圆”；四维篇把下界从 $3.059$ 直接声称到 $4$，属于“把量级从改进变为补全”。矩阵里这两格的差异，对应两种性质完全不同的验证负担（第 9 章）。

\textbf{线索三：跨门派密度上升，验证状态断档。}前五行方法的工具几乎全部来自门内；第六行开始，熵、耦合、钉住、Furstenberg 型估计进入；第七行的外部工具密度最高。与之同步的是验证状态从“已证、同行评审”跌到“未形式化、未评审”。两件事未必有因果关系——批量化生成的工作流天然挤占形式化资源——但它们叠在一起，构成了当前的全部风险来源。

\subsection{矩阵之外的一条观察}

矩阵只列了方法自身。如果把“方法之间的可组合性”也当成一个属性，会看到一条容易被忽略的线索：Guth--Wang--Zahl 的流线化（把 Wang--Zahl 体系改写为凸聚类 + 板条语言）让三维结构第一次可以被“打包”成一块引理；OpenAI 三维篇把这块引理再次加权打包；四维篇因此可以直接调用。方法被反复“打包—复用”是工程成熟的标志，这一点上手稿展示了很高的素养（本文判断，\cfM）。代价同样存在：每多一层打包，验证者就多一层黑盒——核验四维篇的最省力路径，是先核验被它当作黑盒的三维引理链（附录 B）。

\begin{keybox}[可记忆主句]
矩阵最后一行的含义是：结论跨度最大、外部工具最密、验证记录最空——三个“最”同时成立，决定了这份工作此刻只能被称为“声称”。
\end{keybox}
\section{验证缺口与“思想未经严格推演”}

第 4 章记录了“没有 Lean、没有评审”这一事实。本章往前多走一步：把“验证缺口”拆成可以分别追问的三层，并给出一条可执行的核验路线。需要先把分寸立住——本文说“缺口”，说的是证据缺口与理解缺口，不是判定手稿有错。这两件事在数学共同体的日常语用里经常被混为一谈，而它们对研究决策的含义完全不同。

\subsection{三层缺口}

\textbf{第一层：依赖链未验证。}四维篇把自己的启动成本压在了三维姊妹篇的加权板条引理上，后者又立在 Guth--Wang--Zahl 的并集估计上（第 4.3 节）。链条本身没有可疑之处，但意味着验证必须从链条根部做起：先确认三维引理段的推导，再回到四维篇的归约。反过来，如果只核验四维篇第 10 章的收口段，等于在一根未检的承重梁上做局部检修。这类结构性依赖在人类数学里很常见（一篇文章引用另一篇的引理），它之所以在此成为风险，是因为两篇稿件的验证状态同为零，而验证它们的预期难度又都远高于常规文章。

\textbf{第二层：概念未蒸馏。}地平线、两色匹配、水平模型这些自造词汇，目前只在一份手稿里被使用过一次。数学史上，一个概念要“立住”，需要两样东西：别处的独立使用，以及教学意义上的可复述性。以粘性（sticky）为例：Katz--\L aba--Tao 在 2000 年把它作为极端形态的标签提出\upcite{klt2000}，之后二十余年里，它先是变成 Wang--Zahl 的独立定理对象\upcite{wangz2022}，又出现在 Furstenberg 问题、Nikodym 转移与弯曲挂谷的文献中\upcite{wangzakharov2026,gaoliuxi2025,raichoudhuri2026}——概念的跨问题复用，是它被共同体真正消化的实证。地平线会不会走上同一条路，现在无从判断；它甚至连“由人类数学家在讨论班上重新讲一遍”这一步都还没有发生。

\textbf{第三层：推演未受检。}175 页的论证是一个自足的闭环，闭环的正确性只能靠逐段复演来确认。稿件中至少有两类步骤属于“机制最密、最容易藏下隐性假设”的位置：一是比较原理里“导数失配 → 多项式端点映射体积超界”的推导（依赖 Lem\-ma 3.3 型的导数、采样界），二是收口段把局部改进“回填”到原问题的转换（依赖基网格全变差控制）。这两处的共同特点是：如果某条技术性假设过窄，错误不会以“公式矛盾”的方式暴露，而会以“某条不等式在某个子情形下差一个幂次”的方式潜伏——这恰恰是人类审读（以及形式化）最需要耗时的地方。

\begin{keybox}[三层缺口的共同点]
三层缺口指向的都是“证据还没产生”，不是“证据已经反转”。当前对 074 族的正确判断是“未知”，而不是“可疑”——但“未知”在验证成本上等价于一项尚未开工的大型工程，这一点常被讨论忽略。
\end{keybox}

\subsection{“证明”与“理解”是两种成就}

数学共同体对一个结果的接受，实际包含三件事：可核验（有专家能独立复演）、可理解（核心思想能被讲授与改造）、可继承（概念能被其他问题复用）。手稿目前只交付了第四件东西：一份格式完整的文本。

AHM 声明的落点正在“理解”一侧：其结尾呼吁“回到以人类理解为中心的科学愿景”\upcite{ahm2026}。这个立场与“结果是否正确”正交——即使手稿完全正确，理解缺口依然存在。反过来说，也不应把“理解缺口”当作“结果无效”的证据：Newton 时代的许多证明，其可理解性同样是后世几代人才补上的。本文的立场是把两笔账分开记：正确性账（技术层，第 9.4 节的路线图服务于它）与理解账（概念层，需要教学、综述与再表述工作，第 10 章第 3 轨服务于它）。

\subsection{Lean 形式化能覆盖什么、覆盖不了什么}

形式化在此案中的角色需要精确表述，既不能高估也不能低估。它能做的：把一段论证翻译成机器可检查的项，消除“隐性假设”“手滑的常数”“错误的引用编号”这类机械错误，并且把核验的工作量从“全文重读”压缩为“读形式化脚本”。它做不到的三件事：其一，形式化验证的对象是翻译后的命题，翻译本身的保真度不受保护——近期有专门研究以 Navier--Stokes 的案例说明“自动形式化通过检验”与“自然语言证明正确”之间存在系统性缝隙\upcite{bastounis2026}；其二，形式化不判断结果的新意与概念价值；其三，形式化不产生“理解”。

因此，074 族“未形式化”的含义要这样读：不是“有多一层证书”与“没有证书”的差别，而是“机械核验尚未开工，且该领域已经证明这类开工是可行的”——三维挂谷证明（王虹--Zahl）的 Lean 4 形式化已由中科院 MMAT 与 Project Numina 等合作完成\upcite{numina2026}。同一题材、相近量级，人类侧做到了，074 族没有做。为什么没有做，公开材料没有说明；从工作量与稿件整理状态推测（\cfL），最可能的原因不是“不可能”，而是“没有排进优先级”。

\subsection{可复现核验路线图}

\begin{table}[t]
\centering
\caption{四维稿件的模块化核验路线（编号 M1--M8）}
\label{tab:modules}
\footnotesize
\begin{tabularx}{\textwidth}{@{}L{0.9cm}L{3.4cm}L{1.2cm}Y@{}}
\toprule
\rowcolor{lightbg}\textbf{模块} & \textbf{内容} & \textbf{体量} & \textbf{核验要点与风险} \\
\midrule
M1 & 归约（Theorem 2.17 与第 2 章机件） & 约 20 页 & 决定论证是否启动；核对无测度性假设下每步选择；容量类 $\mathcal{C}(d)$ 两条条件的导出 \\
M2 & 比较原理（第 3 章） & 约 22 页 & 导数/采样界（Lemma 3.3 型）是全部列表匹配的前提；体积超界论证的计数 \\
M3 & 临界路径与改进转换（第 4 章） & 约 24 页 & Proposition 4.11 与 Lemma 4.12 的“回填”；指数约定的全局一致性 \\
M4 & 各向同性（第 5--6 章） & 约 35 页 & 次数恢复（Proposition 6.8）的 Remez/Vandermonde 常数控制；圆锥障碍的完备性 \\
M5 & 保边平面刚性（第 7 章） & 约 14 页 & Proposition 7.6 的“保留原始边”与乘法卷积估计的接口 \\
M6 & 有限窄度与水平模型（第 8--9 章） & 约 37 页 & 临界速率 $2k=1$ 的三腿切换；$2k<1$ 自举的终止性 \\
M7 & 无界窄度与收口（第 10 章） & 约 15 页 & 密度指数漂移在比值中的抵消；最终矛盾的指数差为正 \\
M8 & 几何推论（第 11 章） & 约 7 页 & 转移定理假设（Keleti--M\'ath\'e、Gao--Liu--Xi、Nadjimzadah）逐条对照 \\
\bottomrule
\end{tabularx}
\end{table}

建议顺序与理由：先 M1，再“链条根”（三维篇中加权板条引理所在的一节，结合 Guth--Wang--Zahl 的并集估计原文），然后 M2、M3。这四块合计约 70 页，却是整条链上唯一“不可绕开”的部分——其余模块可以按人力并行。形式化的优先级与之类似：M1 与三维引理段的结构最规整、自足性最强，局部形式化的性价比最高；三区间（M4--M7）体量大且依赖分析机件，放在靠后。

\begin{dirbox}[独立小课题（可作为核验的副产品）]
（i）把三维篇 Lemma 2.3 的“加权凸聚类 → 加权板条界”推导整理成人可读版本，作为核验附记；（ii）对 M1 写一份逐步假设审计（audit），逐行确认没有引入可测性、紧性假设；（iii）对三腿切换与锥障碍构造有限离散模型（如 $\mathbb{F}_p$ 类比或小规模数值实验）做压力测试；（iv）对 Proposition 6.8 尝试交互式证明助手内的局部形式化；（v）对照转移定理假设，核对第 11.4--11.5 节推论的条件表。每条产出独立、可验收，见附录 B 的验收标准建议。
\end{dirbox}

\subsection{为什么“正确”之外还要追问“推演”}

写到这里可以给出本章的收束判断：对 074 族的讨论，最好按“三张账”分开记——正确性账（未开工，但有明确路线）、理解账（未开工，需要概念蒸馏与再表述）、影响账（争议中，取决于共同体如何消化）。把三张账混成一张，就会得到两种常见但都站不住的结论：或者因为“AI 生成”而预判不可信，或者因为“结论漂亮”而预支信任。两种结论的共同错误，是让立场代替证据。

\begin{keybox}[可记忆主句]
手稿最稀缺的东西不是篇幅，而是“第二次使用”：一旦地平线、图卡与比较原理在别人的论文里被独立使用一次，这份工作才算真正进入数学的毛细血管。
\end{keybox}
\section{下一步工作的方向建议}

本章把建议分成四条轨道：验证轨道（最短路径上的核验）、数学推进轨道（问题本身往哪走）、方法回灌轨道（这套装置能否搬家）、工程流程轨道（共同体怎么接住批量产出）。每条建议都锚定手稿明确留下的空白或本文核实过的缺口，并附可行性、成本与收益的量级判断；面向研究生与青年学者的可进入点集中放在最后。

\subsection{轨道一：验证}

\begin{dirbox}[验证轨道]
\textbf{1）模块化审读。}按第 9.4 节的 M1--M8 拆分，先把 M1 与链条根（三维引理段 + Guth--Wang--Zahl 原文）做掉。可行性：高；成本：约 4--6 人月的专家工时（按每模块 20--35 页的密度估，\cfM）；收益：给出“论证是否值得继续投入”的第一道闸门。\\
\textbf{2）局部 Lean 形式化。}优先 M1 与三维引理段；不追求全覆盖，追求“依赖链根部可信”。参照人类侧先例（三维挂谷形式化\upcite{numina2026}）；同时配套做翻译保真度校验，避免自动形式化的语义漂移\upcite{bastounis2026}。可行性：中；成本：高于人工审读、低于全篇形式化；收益：一劳永逸的机械核验层。\\
\textbf{3）离散模型压力测试。}对三腿切换（$2k=1$）与次数恢复（圆锥障碍）构造有限域类比或小规模数值实验，检验机制在退化情形下的稳定性。可行性：高（适合作为学位论文章节）；成本：低；收益：为“机制可信度”提供独立于文本的证据。
\end{dirbox}

\subsection{轨道二：数学推进}

手稿留下的空白是密集的，按“离手稿越近越优先”排序：

\begin{dirbox}[数学推进轨道]
\textbf{1）四维极大函数猜想。}手稿只做了 Hausdorff 维数；二维、三维的极大版本均已解决\upcite{cordoba1977,openai3d2026}，四维极大是同一语法下的下一个未结命题。可行性：低（难度即刻上升一档）；成本：高；收益：若成，将与四维集合定理合成完整图景。\\
\textbf{2）$n\ge 5$ 的框架推广。}手稿目前只给投影下界（$\dim_{\mathrm{H}}\ge 4$）\upcite{openai4d2026}。一个自然的猜想式路线：把“二次图卡”升级为“低次代数曲面图卡”，并把窄度参数 $\ell$ 的区间划分改写为按“拟合次数 vs 时间尺度”的分层——四维里 $\ell=0$ 对应的圆锥障碍，在高维将替换为更广的簇障碍。可行性：未知；成本：极高。\\
\textbf{3）概念蒸馏与再表述。}把地平线、两色匹配、水平模型写成独立的、可讲授的材料（讨论班讲义、教程），检验这些概念是否独立成立、是否可以被改造后用于同类问题（如投影理论、Radon 型变换）。可行性：高；成本：低；收益：直接服务“理解账”，且这类工作天然适合作为综述与教科书的章节。\\
\textbf{4）定量联动。}三维篇已给出极大估计的 $\delta^{-\varepsilon}$ 版本，与 Fourier 限制、Bochner--Riesz、局部光滑的联动在三维可以立即着手；球面限制旁路篇（family 077）与挂谷主线的关系值得专文对照\upcite{openaielliptic2026,openai3d2026,gwz2026}。
\end{dirbox}

\subsection{轨道三：方法回灌}

\begin{dirbox}[方法回灌轨道]
\textbf{1）地平线框架的移植。}“局部拟合精度 $\leftrightarrow$ 时间代价”的记账方式不依赖四维几何本身；把同一装置移植到：限制型算子的波包聚集问题、Furstenberg 型投影问题、以及 Kakeya--Nikodym 的极值分析，是低风险高观测价值的试探。可行性：中；成本：中。\\
\textbf{2）熵账本与平面刷的融合。}门内的平面性/颗粒性机制（平面刷）与借入的熵/耦合机制目前是两套语言；若能用一套统一表述覆盖两者，将显著降低后续论文的阅读成本。可行性：中；成本：中；收益：为下一代工作提供公共语言。\\
\textbf{3）把“打包—复用”变成常规。}Guth--Wang--Zahl 流线化 → OpenAI 三维加权化 → 四维复用，这条打包链展示了可组合性红利；建议将“引理接口的稳定性”当作独立指标来维护（例如在形式化库里固定接口），这也回应了第 9 章的验证瓶颈。可行性：中高；成本：中。
\end{dirbox}

\subsection{轨道四：工程流程}

\begin{dirbox}[工程流程轨道]
\textbf{1）分诊协议。}对批量产出建立公开的分诊表：按“重要性 × 可验证性 × 人力预算”排序，先做依赖链根部与高复用潜力的结果。本综述的附录 B 即是一份微型样例。\\
\textbf{2）形式化覆盖率作为发布契约。}仓库已按版本管理修订（撤稿、补丁、形式化增量）\upcite{openairepo2026}；社区侧的合理诉求是把“顶层结果形式化覆盖率”与“独立审读状态”纳入发布标准——AHM 声明要求的“停止合作”与 OpenAI 主张的“继续对话”之间，至少存在一个技术性共识区：验证基础设施的共建\upcite{ahm2026,openairelease2026,tao2026}。\\
\textbf{3）评价指标的补齐。}当前对这类发布，唯一可公开核对的指标是形式化覆盖率（300/719 $\approx$ 42\%）；真正决定长期影响的“独立审读周期”与“概念复用率”尚无任何制度性跟踪。建议以 074 族为试点建立这两项记录。
\end{dirbox}

\subsection{面向研究生的可进入点}

\begin{table}[t]
\centering
\caption{具体可进入的小课题与验收标准}
\label{tab:entrypoints}
\footnotesize
\begin{tabularx}{\textwidth}{@{}L{0.6cm}L{4.6cm}Y@{}}
\toprule
\rowcolor{lightbg}\textbf{编号} & \textbf{题目} & \textbf{验收标准} \\
\midrule
E1 & 三维篇加权板条引理的独立复演（含与 Guth--Wang--Zahl 并集估计的接口） & 交付一份自足推导（10--20 页），专家复核无缺环；并记录所有用到的假设 \\
E2 & 四维篇归约（M1）假设审计 & 逐步清单：每一步的假设、结论、隐含常数依赖；确认无测度性假设混入 \\
E3 & 三腿切换的离散模型实验 & 可复现代码 + 报告：在 $\mathbb{F}_p$ 类比或数值模型中观察/反驳机制；记录失败情形 \\
E4 & Proposition 6.8 的交互式形式化 & 可编译的局部形式化脚本（允许带未证洞），配翻译对照表 \\
E5 & 推论章转移定理条件核对（第 11.4--11.5 节） & 条件对照表：每条转移定理的假设 vs 手稿使用处，逐条标绿/标红 \\
\bottomrule
\end{tabularx}
\end{table}

E1 与 E2 是“小投入快反馈”的一对：不要求核验者同意手稿的全部机制，只要求把链条根部的推导独立走一遍——数学上干净，方法上保守，产出立刻可被共同体使用。E3 适合有计算背景的学生；E4 适合形式化方向；E5 适合偏好文献与结构分析的工作。

\begin{keybox}[可记忆主句]
这份手稿给学界留下的最有价值的东西，可能不是那个待验的四维结论，而是一张已经画好的攻击地图：谁先按图索骥核验链条根部，谁就掌握了后续所有工作的闸门。
\end{keybox}
\section{修正后的判断}

导言里悬置的问题，到这里可以给出修正后的判断了。所有判断按本文的置信度分级标注，并附与第 2 至 10 章的对应依据。

\textbf{判断一：技术层面，若手稿成立，这是一次问题层面的范式级推进。}\cfH{}三窄度区间 + 地平线的组合，是对四维特有障碍（低次代数曲面上的聚集）的第一份系统性回应；它放弃了门内“形态穷举”的路线，改用参数分层与极值改进，方法的完整度从第 5 章的骨架重建来看是自洽的。这里的“范式级”限定在问题内部——它替换的是解决四维挂谷的组织方式，而不是数学思想的坐标系；三层分解见判断七。但“若成立”三个字承载的分量很大——本判断的可信度取决于验证状态，而不是文本的精致程度。

\textbf{判断二：谱系层面，这是挂谷门派有史以来距离最远的一次跨门派借入。}\cfH{}熵账本、概率耦合、实代数几何纪律与乘法卷积 Fourier 衰减在门内文献中零命中（第 7.3 节词频表），这一事实不依赖于对手稿正确性的判断。历史经验（加法组合、多项式方法的消化各用了一代人）提示：借入距离越远，理解成本越高，而理解成本正是当下最大的隐性支出。

\textbf{判断三：对数学共同体，验证瓶颈已取代解题难度成为主要矛盾。}\cfH{}一个可以从集群侧无限提速的生成端，对接一个只能线性扩展的验证端，“生成快、验证慢”是结构性不对称；074 族是这个不对称的第一次标志性显影——272 页、零形式化、零审评，而全世界的合格审读者只有两位数。

\textbf{判断四：形式化是部分答案，且只覆盖机械层。}\cfM{}它在依赖链核验上不可替代，人类侧的三维挂谷形式化已经证明了可行性；但它既不能担保自然语言层与形式层的语义一致，也不能产生概念理解。把它当作验证的全部，或当作可以忽略的余事，都会误判成本结构。

\textbf{判断五：对读者与青年学者的使用建议。}\cfM{}把这两篇稿件当作一份“待核验的研究计划书 + 工程样机”，而不是定理库中的新条目：先读第 5 章的骨架、附录 B 的清单，再决定投入方向；核验的副产品（E1--E5 类小课题）本身具有独立发表价值。任何“相信”或“不信”的整体表态，都不如先完成链条根部的一次复演。

\textbf{判断六：对“AI 辅助数学”这一宏观问题的姿态。}\cfM{}本案例既支持乐观叙事（方法论创新密度确实高于同期人类论文的平均值），也支持悲观叙事（验证与理解的欠账被同步放大）。两种叙事共享一个事实基础：数学共同体的核心资产不是结论库，而是“理解—复用—再生成”的循环；批量产出对这个循环的冲击，远大于对结论库的冲击——这也正是 AHM 声明与 OpenAI 发布口径之间争议的真正落点\upcite{ahm2026,openairelease2026,tao2026}。

\textbf{判断七（“开创还是闭包”的分层定位）。}\cfM{}把“是否发明了新数学思想”拆成三层来问，答案才站得住。\emph{命题层}：$\dim_{\mathrm{H}}K=4$ 若成立，属于发现——它陈述的是对象本来具有的性质，无所谓发明，只有真伪。\emph{装置层}：图卡、地平线、比较原理、次数恢复、双帧实验、保边回填这些对象与机制属于发明——它们在门内文献中零命中（第 7.3 节词频证据），但每一项都有可辨认的亲缘（多项式 Remez、划分、极值反证、熵、耦合），是“在既有语言内新造机件”。\emph{范式层}：极值反证加局部拟合加逐尺度归纳的大框架延续自门内传统；本次真正带思想性质的新意，是把“局部理解所需付出的代价”形式化为可极值化的资源（地平线），并首次把信息论、分形几何与实代数几何三套语言组装进同一套证明纪律。据此，本文的定位是：\textbf{它不是“现有知识方法的凸闭包推论”}——该闭包在 $3.059$ 与 $4$ 之间停摆了二十余年，且证明需要新增对象；\textbf{也够不上“革命性的数学思想创新”}——它没有像多项式方法或熵方法当年那样重写整个领域的问题意识。最准确的刻度是：\emph{继承框架、发明装置、跨门派综合的破局性工作}——对问题本身是决定性跃迁，对数学思想是高水平综合，对 AI 参与数学的工具箱则是第一次“全域架构级”产出。以上以手稿成立为前提；若最终被修正，装置层的评估需按存留的有效部件重新计算。三层结论的技术底稿分别为：装置的精确定义与推导细节（附录 C.1）、与传统极值反证的形式化对比（附录 C.2，含 $H(E)$ 与改进引理的公式化写法）、三套外门派语言的映射表（附录 C.3）。

回到第 2 章的问题本身：方向完备性与空间压缩的张力，在四维以上还没有被完全驯服——即便手稿正确，它给出的也只是 Hausdorff 维数这一条战线上的胜利，极大函数、更高维、正测度之外的精细结构仍在原位。数学的进展向来如此：把一条防线推到理论上限，剩下的防线立刻变成新的前沿。这份稿件若被验证，它值得的位置是“四维战场的转折点”；若被修改或证伪，它同样已经证明了一件事——挂谷问题最前沿的攻击面，已经可以被机器系统性地测绘出来。

\begin{keybox}[可记忆主句]
对这份手稿，正确的姿态既不是膜拜也不是鄙视，而是把它放进验证流水线：先复演链条根部，再谈四维挂谷的前途；在证据到来之前，未知就是未知。
\end{keybox}

\appendix
\section{记号与依赖图}

\subsection{记号速查表}

\begin{table}[ht]
\centering
\caption{本文与两篇稿件共用的主要记号}
\label{tab:notation}
\footnotesize
\begin{tabularx}{\textwidth}{@{}L{3.5cm}L{5.6cm}Y@{}}
\toprule
\rowcolor{lightbg}\textbf{记号} & \textbf{含义} & \textbf{出处/备注} \\
\midrule
$K,\ E$ & $\mathbb{R}^n$ 中的挂谷集（含每个方向的单位线段） & 第 2 章；两篇稿件通用 \\
$\dim_{\mathrm{H}},\ \underline{\dim}_{\mathrm{M}},\ \overline{\dim}_{\mathrm{M}}$ & Hausdorff、下/上 Minkowski 维数 & 第 2 章式 (2.1) \\
$T_\delta(a,e)$ & 中心 $a$、方向 $e$、长 $1$、半径 $\delta$ 的细管 & 四维篇 §1.1 \\
$K_\delta$ & 挂谷极大算子（对细管内平均取方向上的上确界） & 两篇稿件通用 \\
$Y(T),\ \lambda$ & 细管的阴影及其占管体积的比例 & 三维篇 (1.1) \\
$N_0,\ r,\ L$ & 原问题分辨率、观察深度、终端深度 & 四维篇 §2.1 \\
$p_r,\ \bigl[w\bigr]_{a_r/B}$ & 深度 $r$ 的观察；隐藏坐标 $w$ 的分箱 & 四维篇式 (2.3) \\
$n(r)=N_0^{-f(r)}$ & 典型入射密度；$f$ 非降、1-Lipschitz & 四维篇 §2.1 \\
$d<1$ & 容量类 $\mathcal{C}(d)$ 的密度退化率：$f(L)-f(r)\le d(L-r)$ & 四维篇 Definition 2.6 \\
$h,\ H(E)$ & 图卡的时间代价；地平线（正时间代价下确界） & 四维篇 §1.2、§2 \\
$k,\ \ell$ & 临界路径的时间速率；最优窄度（$\ell=0$ 各向同性，$0<\ell<\infty$ 有限窄度，$\ell=\infty$ 无界窄度） & 四维篇 §1.2、§4 \\
版本 1/版本 2 & 精确问题的两个模型：标量模型（$j=1$）与四维模型（$j=2$） & 四维篇 Definition 2.1 \\
$Y'=-Z+tv$ & 有限窄度区间的水平模型；盒型水平尺度 $H$、横向 $H^2$ & 四维篇 §1.2、§8 \\
$h(p,z)$ & 三维篇的临界指数（“极大定理成立”等价于 $h=0$） & 三维篇 Definition 2.3 \\
$F(s)=\beta(s-\tau)_{+}$ & 取等配置的规范时间剖面（平段 + 恒定亏空斜率） & 三维篇 §1.2、§5 \\
$X,Y,U,V$ & 双帧实验的四坐标（法向位置/切向位置/法向速度/切向速度） & 三维篇 §1.2、§7 \\
$\Delta t\,\Delta\vartheta$ & 双帧的乘积斜率（钉住估计的控制对象） & 三维篇 §1.2、§9 \\
\bottomrule
\end{tabularx}
\end{table}

\subsection{依赖图}

图 \ref{fig:deplink} 把第 2--5 章涉及的主要依赖画成一张图：主干是“历史下界演进 → 三维全解 → OpenAI 三维篇 → OpenAI 四维篇 → 三区间改进”的纵向链条；右侧虚线是四维篇从其他门派借入的工具（第 7 章逐一分析）；右下角的虚线框是三维篇与球面 Fourier 限制篇之间的旁路。

\begin{figure}[t]
\centering
\begin{tikzpicture}[
  node distance = 5.6mm,
  box/.style = {draw = lightline, line width = 0.4pt, fill = lightbg, rounded corners = 2pt,
    align = center, font = \footnotesize, inner sep = 3.2pt, text width = 76mm},
  sidebox/.style = {draw = lightline, line width = 0.4pt, fill = lightbg, rounded corners = 2pt,
    align = center, font = \scriptsize, inner sep = 3pt, text width = 46mm},
  flow/.style = {-{Stealth[length = 2.1mm]}, draw = navy, line width = 0.7pt},
  side/.style = {-{Stealth[length = 2.1mm]}, draw = sky, dashed, line width = 0.6pt}
]
\node[box] (h1) {Wolff 1995 发刷：$\mathbb{R}^4$ 下界 $3$；多项式 Wolff 公理时代：$3+\frac{1}{40}$；平面刷：$3.059$};
\node[box, below = of h1] (h2) {Wang--Zahl 2022--2025（粘性 $\to$ 三维全解）；Guth--Wang--Zahl 2026 流线化（凸聚类并集估计）};
\node[box, below = of h2] (o3) {OpenAI 三维篇：$\lambda^{3}$ 极大估计；加权全时段板条界（其 Lemma 2.3）};
\node[box, below = of o3] (o4) {OpenAI 四维篇 Theorem 2.17：正地平线 $H(E)>0$ 的加权问题};
\node[box, below = of o4] (rr) {三区间改进（各向同性 / 有限窄度 / 无界窄度）$\Rightarrow$ 与极值性矛盾 $\Rightarrow \dim_{\mathrm{H}}K=4$};
\draw[flow] (h1) -- (h2);
\draw[flow] (h2) -- (o3);
\draw[flow] (o3) -- (o4);
\draw[flow] (o4) -- (rr);

\node[sidebox, right = 16mm of o4] (s1) {多元 Kakeya（行列式形式）；Wang--Zahl 限制三元组点积定理};
\node[sidebox, below = 6mm of s1] (s2) {Orponen--de Saxc\'e--Shmerkin 乘法卷积 Fourier 衰减；图版本 BSG 与和集演算；实代数几何与 Bernstein 类};
\node[sidebox, below = 6mm of s2] (s3) {三维篇旁路：球面 Fourier 限制篇（独立路径，虚线）};
\draw[side] (s1.west) -- (o4.east);
\draw[side] (s2.west) -- (o4.east);
\draw[side] (s3.west) -- (o3.east);
\end{tikzpicture}
\caption{核心依赖图。实线为主干依赖，虚线为跨门派工具输入（第 7 章）与旁路（球面 Fourier 限制篇）。}
\label{fig:deplink}
\end{figure}
\section{核验优先级清单}

本附录把第 9.4 节的路线图落成可打勾的清单，供核验小组直接使用。每条给出：核验对象、依赖、建议投入量级、完成标志（DoD，Definition of Done）与常见陷阱。

\subsection{第一优先级：链条根部（建议先行）}

\begin{table}[ht]
\centering
\caption{第一优先级核验清单}
\label{tab:check1}
\footnotesize
\begin{tabularx}{\textwidth}{@{}L{2.4cm}L{2.4cm}Y@{}}
\toprule
\rowcolor{lightbg}\textbf{对象} & \textbf{依赖} & \textbf{完成标志与陷阱} \\
\midrule
三维篇加权板条引理段（含加权凸聚类推导） & Guth--Wang--Zahl 并集估计原文\upcite{gwz2026} & DoD：自足推导一份，专家复核无缺环；陷阱：凸聚类条件与阴影密度阈值之间的换算是全链最容易“差一个幂”的位置 \\
四维篇 M1（Theorem 2.17 归约） & 上一条 + 四维篇第 2 章 & DoD：逐步假设审计表（E2）；陷阱：外层测度可数次可加性选段处的可测性处理、容量类两条条件的导出顺序 \\
三维篇极大估计的主定理路线（$h(p,z)=0$） & 三维篇 §2--§10 & DoD：把“临界指数为正 $\Rightarrow$ 取等配置”的机件复演一遍；陷阱：规范剖面 $F(s)=\beta(s-\tau)_{+}$ 的最小推迟论证 \\
\bottomrule
\end{tabularx}
\end{table}

\subsection{第二优先级：机制密度最高的两处}

\begin{table}[ht]
\centering
\caption{第二优先级核验清单}
\label{tab:check2}
\footnotesize
\begin{tabularx}{\textwidth}{@{}L{2.4cm}L{2.4cm}Y@{}}
\toprule
\rowcolor{lightbg}\textbf{对象} & \textbf{依赖} & \textbf{完成标志与陷阱} \\
\midrule
四维篇 M2（比较原理，第 3 章） & M1 & DoD：列表匹配的体积论证逐步复核；陷阱：导数/采样界的边界情形（列表“频繁取值一致”的定量含义） \\
四维篇 M3（Proposition 4.11 与 Lemma 4.12） & M1、M2 & DoD：回填步骤的全变差估计独立推导；陷阱：原律恢复时“不向列表规则暴露未观测标签”的约定 \\
\bottomrule
\end{tabularx}
\end{table}

\subsection{第三优先级：三区间与收口（可并行）}

M4--M7 按第 9.4 节的表 \ref{tab:modules} 执行；配套独立课题 E1--E5 的验收标准见第 10 章的表 \ref{tab:entrypoints}。三区间内部的建议顺序：各向同性（M4，含次数恢复）$\to$ 有限窄度（M6，含水平模型与临界速率）$\to$ 无界窄度（M7，含收口）。

\subsection{通用核验规范（建议）}

\begin{dirbox}[核验工作规范]
（i）每一步的结论必须写成“假设 → 断言 → 常数依赖”的三元组，禁止转述性省略；（ii）任何“明显地”“容易看出”手稿用语，一律展开为显式推导；（iii）所有指数比较必须给出正间隙的数值下界（手稿的全局约定是在分辨率趋于无穷前保持固定正间隙，核验时沿用\upcite{openai3d2026,openai4d2026}）；（iv）每日核验记录以“可复核的最小反例搜索”收尾：对小规模离散模型做一轮反例扫描；（v）发现缺口时按手稿的修订协议格式登记（对象、位置、缺口类型、最小修复猜想），便于与发布方版本管理对接\upcite{openairepo2026}。
\end{dirbox}

\begin{cautbox}[三张账分开记]
正确性账：按本清单走；理解账：概念蒸馏与再表述（第 10 章轨道二第 3 条）；影响账：独立审读周期与概念复用率的长期跟踪（第 10 章轨道四第 3 条）。三张账混记，是当前讨论最大的方法论混乱之源。
\end{cautbox}
\section{核心装置：精确定义、极值原理与语言映射}

本附录为第 5、7、11 章提供精确支撑，分三部分：C.1 转录四维篇关键装置的定义与命题形态（含编号定位，便于逐条核对）；C.2 把“地平线”与传统极值反证做形式化对比；C.3 给出三套外门派语言的映射表。需要先声明一句：以下定义与命题均属\textbf{手稿声称的结构}，本文转录它们是为了说明“装置长什么样”，不构成对其正确性的背书。

\subsection{装置的精确定义（依四维篇编号）}

\subsubsection*{图卡、真图卡与已知图集（Definition 2.8）}

图卡系统在深度 $r\in(0,L]$ 上把一个公共二分时间划分成区间 $I$，长度
\begin{equation}
q=N_0^{-h+o(1)},\qquad h\ge 0,
\end{equation}
在每个世界与每个区间内把不相交的轨迹集 $A_c$ 分配给图卡标签 $c$，选定入射的总质量至少 $K_0^{-1}$。每张图卡配一组几何数据：仿射空间中心 $y_c(t)$ 与可逆矩阵 $D_c$，满足
\begin{equation}
\|D_c\|\le qK_0,\qquad J_c=|\det D_c|,\qquad \sup_{t\in I}\bigl|D_c^{-1}(y(t)-y_c(t))\bigr|\le K_0\quad (l\in A_c),
\end{equation}
并配一个二次多项式测试 $P$（“测试保持在允许厚度内”的精确含义）：
\begin{equation}
\begin{aligned}
&\sup_{t\in I}|P(t,y(t),w(t))|\le K_0a_r,\qquad
\sup_{t\in I}|\partial_wP|\le K_0B,\\
&|\partial_wP|\ge B/K_0\ \text{（在选定入射上）},\qquad
|\partial_{ww}P|\le K_0B^2/a_r .
\end{aligned}
\end{equation}
两条质量分别记录：分配轨迹测度 $\nu_c$ 与选定入射测度 $m_c$。\textbf{真图卡}的条件是
\begin{equation}
\frac{m_c}{q}\;\ge\;K_0^{-1}\,n(r)\,J_c,
\label{eq:true}
\end{equation}
即“单位时间内的入射质量”不低于典型密度乘空间体积代理；系统全由真图卡组成时称为真系统。\textbf{已知图集}是标签可由观察（按 $p_L$）从一份确定性列表读出的图卡系统，它保留分配、覆盖与温度性（temperedness）要求，但不必满足式 \eqref{eq:true}。三个量的分工在这里已经清楚：$q$ 是时间分辨率，$m_c$ 是覆盖证据，而式 \eqref{eq:true} 把“局部理解”变成了可加比较的量化陈述——这正是它区别于“几何形态假设”（如粘性）的地方。

\subsubsection*{地平线（§1.2 与第 2 章）}

把区间长度写成 $N_0^{-h+o(1)}$ 时的 $h$ 称为\textbf{时间代价}。记 $\mathcal{T}(E)$ 为加权问题 $E$ 的全部真终端图卡系统，$m(\Sigma)$ 为系统 $\Sigma$ 的覆盖量（$N_0^{-o(1)}$ 量级称实质覆盖），则\textbf{地平线}定义为
\begin{equation}
H(E)\;=\;\inf\bigl\{\,h(\Sigma)\;:\;\Sigma\in\mathcal{T}(E),\ m(\Sigma)\ge N_0^{-o(1)}\,\bigr\}.
\label{eq:horizon}
\end{equation}
归约定理（Theorem 2.17）的形态是：若 $K\subset\mathbb{R}^4$ 含每个方向的单位线段且 $\dim_{\mathrm{H}}K<4$，则存在 $d<1$ 与属于容量类 $\mathcal{C}(d)$ 的温和齐次精确问题 $E$，使 $H(E)>0$。注意这里的“极值化”是双层结构：外层是 $H(E)$ 对系统的下确界；内层是容量类对密度退化率 $d<1$ 的约束——$d<1$ 恰是比较原理（C.1 下一小节）得以启动的前提。

\subsubsection*{代数叶片与两色匹配（Lemma 3.1--3.3）}

\textbf{代数叶片}指二次测试的局部图像所给出的解析分支：在某条比较路径的邻域内，测试方程的解是函数变量上的仿射函数、线性方程的根或二次方程的简单根，这些分支沿路径可解析延拓。手稿给出的关键定量事实（Lemma 3.2，\emph{延拓的多项式代价}）是：在 $A^{-1}$ 邻域内、端点比较球（半径 $c/A$）上，延拓函数的范数双向比较损失至多 $A^{C}$（$C$ 与方程系数无关），固定阶导数由这些范数控制；该一致界经“半代数描述 + 量词消去”得到——即\emph{常数纪律本身被写成可验证的代数事实}。

两色匹配（Lemma 3.3）是全部列表匹配的定量前提，其否定式陈述为：固定 $D\ge1$、整数 $J>D$、$d_0<1$，不存在 $N\to\infty$ 的序列，同时满足
\begin{equation}
\#\{\text{每个列表}\}\le N^{d+o(1)}\ (d\le d_0),\qquad
\|f\|_{C^{J+1}}\le K=N^{o(1)},\qquad
\text{边际约束}\ \le K\,\pi_{\kappa,i}\,dt,
\end{equation}
\begin{equation}
|f_{1,i_1}(t)-f_{2,i_2}(t)|\le KN^{-1}
\quad\text{且}\quad
\max_{1\le j\le D}\bigl|f^{(j)}_{1,i_1}(t)-f^{(j)}_{2,i_2}(t)\bigr|\ge K^{-1}
\quad\text{（在保留下来的事件上）}.
\end{equation}
直白地说：值接近（$N^{-1}$ 精度）与导数分离（$K^{-1}$ 下限）在同一采样律下不可长期共存，只要列表规模小于任何完整幂（$d<1$）。证明机制是“交替路径 + 端点映射的 Jacobian”：沿两列表交替取样构成多项式端点映射，其像迫使记录数超过列表规模允许的上限；手稿自述这一用法与 Christ 的 incidence 算子细化法（method of refinements）类比\upcite{openai4d2026}。

\subsubsection*{改进的定量定义与转换（Definition 4.10、Proposition 4.11、Lemma 4.6）}

\textbf{时间改进候选}：在准备好的切父块中，对厚度 $a=N^{-s}$、时间长度 $q_s$、轨迹质量 $\nu_Q$ 的路径包 $Q$，一个增益 $c>0$ 的时间改进候选是宽度 $a'=aN^{-c}$ 上的二次测试，满足
\begin{equation}
\begin{aligned}
&|P|\le Ka',\qquad |P_w|\le K,\qquad |P_{ww}|\le K/a',\qquad
|P_w|\ge K^{-1}\ \text{（计数入射上）},\\
&m\ge K^{-1}N^{-dc}\nu_Q .
\end{aligned}
\end{equation}
当 $0<\ell<\infty$ 时，记 $g=N^{-\ell s}$，\textbf{条带改进候选}是对固定行 $R$ 要求 $|Y'-R\cdot(1,v)|\le KgN^{-c}$，且作用在质量至少 $K^{-1}\nu_Q$ 的轨迹片上。两句解释性原文值得直译：时间改进候选“在不改变时间块的前提下改进隐藏坐标拟合”，条带改进候选“改进横向速度预测”——改进不是模糊的“更精细”，而是两个具体自由度（隐藏坐标拟合、横向速度）上的定量收益。

\textbf{转换与排除}：Proposition 4.11 把“每个实质残余上都存在候选”升级为原分辨率下的改进输出（含 Lemma 4.12 的列表回填）；Lemma 4.6 则是一份“有限样本排除”清单：它排除的那类原分辨率输出不允许存在。两者相扣，候选假设即不可能——这是四维证明的最终矛盾形态。

\subsubsection*{极值临界路径与三区间的收口（Proposition 4.3、Theorem 6.1、Proposition 10.1）}

Proposition 4.3（\emph{临界路径}）把“极大序列”规范化：长度为一的精确问题的每个最大化序列满足相对密度剖面线性化（斜率 $d^{*}$），且对每个有限有序深度网格存在嵌套真系统，其累积地平线趋于 $kr$、按自身长度单位的分段地平线趋于 $k$——$k$ 即“时间速率”。另一参数 $\ell$ 是同一极值安排下“较小空间轴的优化窄度”。三区间的收口分别落在：各向同性（$\ell=0$）由 Theorem 6.1 给出改进——在每个实质残余中找到厚度 $a=N^{-s}$、长度 $h\sim_{\log,N}a^{k}$ 的包，以及厚度
\begin{equation}
a_2=a\,h^{e_2}\quad (e_2>0),\qquad
m\ge \frac{h\nu_Q}{K}\Bigl(\frac{a_2}{a}\Bigr)^{d},
\end{equation}
的时间改进候选；有限窄度（$0<\ell<\infty$）由第 8--9 章的标量投影、水平模型与临界速率 $2k=1$ 处理；无界窄度（$\ell=\infty$）由 Proposition 10.1 排除——在第二模型的极值设置下全局优化窄度 $\ell$ 有限。

\begin{keybox}[为什么这些构成“发明”而非“换名”]
以三条判据自查，四件装置都过关（本文判断，\cfM）：\textbf{（i）对象此前不存在}——图卡、地平线、两色匹配、改进候选在门内文献中零命中（第 7.3 节词频证据）；\textbf{（ii）角色不可替代}——真性条件 \eqref{eq:true} 堵住“只拟合例外族”的漏洞，两色匹配供给跨尺度延拓的定量前提，改进定义 \eqref{eq:horizon} 之后的候选--转换对才是矛盾引擎；\textbf{（iii）非旧物换名}——“真图卡”是拟合质量判据，与“粘性”这类几何形态假设在逻辑角色上不同：前者可被逐块量化比较，后者是对象的整体性质。判据（iii）也给出反面对照：如果某装置只是把已知对象改换名称，它应当能直接继承旧对象的既有定理；而这四件装置都需要伴随的定量引理（式 \eqref{eq:true}、Lemma 3.2、Lemma 3.3、Lemma 4.6）才能工作。
\end{keybox}

\subsection{地平线与传统极值反证：形式化对比}

\subsubsection*{两种范式的形式骨架}

\textbf{传统型（取等分析作引擎）}——同一批稿件中的三维篇就是现成样本：假设临界指数 $h(p,z)>0$，由失败导出\textbf{取等序列}（equality sequences）；取等配置的几何被逼进一条规范时间剖面 $F(s)=\beta(s-\tau)_{+}$ 与“同一指标双帧”的投影约束实验；最后按形态逐类排除，与熵需求矛盾\upcite{openai3d2026}。其形式骨架是：
\begin{equation}
\text{(T1)}\ \text{估计失败}\; \Longrightarrow\; \exists\ \text{取等序列};\qquad
\text{(T2)}\ \text{取等对象落在有限个形态类};\qquad
\text{(T3)}\ \text{逐类矛盾}.
\end{equation}
极值对象是\emph{几何配置}（管族、阴影、棱柱、板条）；(T2) 的可行性依赖紧性与极限结构，把“列不出形态”视为主要风险。

\textbf{地平线型（以代价下降为引擎）}——四维篇的骨架是：
\begin{equation}
\text{(H1)}\quad \dim_{\mathrm{H}}K<4\ \Longrightarrow\ \exists E:\ H(E)>0,
\end{equation}
\begin{equation}
\text{(H2)}\quad \text{极大路径规范化：相对剖面线性、时间速率 } k \text{、优化窄度 } \ell ,
\end{equation}
\begin{equation}
\begin{aligned}
&\text{(H3)}\quad \text{在每个实质残余上存在改进候选}\\
&\qquad\Longrightarrow\ \text{(Prop 4.11) 原分辨率改进输出}\ \Longrightarrow\ \text{(Lemma 4.6) 排除}.
\end{aligned}
\end{equation}
极值对象从\emph{配置}换成了\emph{描述系统} $(\Sigma,h,m)$：被极值化的是“把拟合做到终端厚度”的时间代价 \eqref{eq:horizon}，矛盾项不是“某类形态不可能”，而是“下确界的任何逼近都可被严格改进”。

\subsubsection*{三条本质区别}

\textbf{对象之变：从配置到“描述--代价对”。}传统型的变量是研究对象本身（管族如何重叠、在哪里聚集）；地平线型的变量是关于对象的局部知识（图卡系统对轨迹的拟合质量与覆盖）。把知识状态而非几何状态当作极值对象，是这套机制与门内一切既有反证的第一个断点。

\textbf{引擎之变：从分类排除到普适下降。}传统型的矛盾靠“形态类有限且逐类可杀”；地平线型不需要形态清单，它证明的是一句全称命题——\emph{任何}近乎极值的系统都存在下降方向（时间代价降低，或固定时间下的窄度增益）。分类式论证怕“漏掉新形态”（四维聚集形态可能列不尽），下降式论证不怕：它只需在每个残片制造一次收益，收益本身就是矛盾。

\textbf{前提之变：从紧性到一致性比较。}取等分析需要序列收敛到可分析的极限对象，因此依赖紧性、正则性与极限理论；下降论证不构造极限对象，它依赖的是“局部候选 -> 原问题改进”的转换（Proposition 4.11 与 Lemma 4.12），其技术前提是列表一致性（两色匹配、全变差回填）与指数账目的一致性。用一致性与记账替代紧性，是这套机制可以绕开“四维没有可用紧类”这一死结的原因（本文判断，\cfM）。

\begin{table}[t]
\centering
\caption{两种极值反证范式的对照}
\label{tab:extremal}
\footnotesize
\begin{tabularx}{\textwidth}{@{}L{2.7cm}L{5.3cm}Y@{}}
\toprule
\rowcolor{lightbg}\textbf{维度} & \textbf{传统型（取等分析）} & \textbf{地平线型（代价下降）} \\
\midrule
极值对象 & 几何配置（管族/阴影/板条） & 描述系统：图卡系统与其覆盖、时间代价 $h$ \\
反证输入 & 估计失败或临界指数为正 & 维数亏空 $\Rightarrow H(E)>0$ \\
存在性原理 & 失败 $\Rightarrow$ 取等序列（需要极限对象） & 亏空 $\Rightarrow$ 正代价（下确界下界；不构造极限对象） \\
矛盾引擎 & 形态分类 $\to$ 逐类排除 & 普适下降：每个实质残余存在改进候选 $\to$ 转换 $\to$ 排除 \\
技术前提 & 紧性、正则性、极限理论 & 两色匹配、全变差回填、指数账目一致性 \\
结论形态 & “不存在取等对象” $\Rightarrow$ 估计成立 & “下确界不可被逼近” $\Rightarrow$ 正下界不可能 \\
\bottomrule
\end{tabularx}
\end{table}

\begin{cautbox}[诚实的边界]
三条区别描述的是机制结构，不是绝对的谱系断绝。地平线型与“最小反例归纳法”“变分直接方法”都有亲缘：它同样以“极值不可达”收尾，也借用了“改善方向”的直觉。真正的差别在于：它没有把极值对象嵌入紧空间，也没有欧拉--拉格朗日方程，下降发生在离散的描述系统上，由组合与指数账目驱动。因此本文把它定位为\textbf{思想性新意}而不是新范式（\cfM）：它给“极值反证”这一古老工具增加了一个新维度——“关于知识的代价”这一维度。另一个值得记录的旁证是：在同一批稿件内部，三维篇仍走传统取等路线（$F(s)=\beta(s-\tau)_{+}$ 剖面），四维篇才切换为地平线——两代机制的对照在同一系统内部完成，说明这不是评审语言的复制，而是对四维障碍的定向回应（本文判断，\cfM）。
\end{cautbox}

\subsection{三套外门派语言的映射表}

判断七所称“三套语言”指信息论、分形几何（熵与投影）与实代数几何。表 \ref{tab:langmap} 逐条给出它们的挂谷侧对象与落点命题；三条备注说明分寸：映射不是名词替换，每条都在手稿中被重新证明为定量形式（均有专节）；第四项小借入（极值组合的图版 BSG 与和集演算）见 §7，不纳入三语言框架；钉住（pinning）与乘法卷积是分形几何语言的内部构件，故并入第二行组。

\begin{table}[t]
\centering
\caption{三套外门派语言与四维挂谷问题的映射}
\label{tab:langmap}
\footnotesize
\begin{tabularx}{\textwidth}{@{}L{2.5cm}L{4.4cm}Y@{}}
\toprule
\rowcolor{lightbg}\textbf{语言与构件} & \textbf{挂谷侧对象} & \textbf{映射与落点命题} \\
\midrule
信息论：条件熵、互信息 & 观察的鉴别力（隐藏坐标分箱携带的信息量） & 三次、四次项失配的检测与“熵窗口”选择（四维篇 §5，Lemma 5.3 型熵界）；参数分辨率提升的自举（§9） \\
信息论：联合律与乘积律的耦合比较 & 同一指标的两次观测（双帧） & 两事件共享指标时，实际联合律与乘积律的偏差逼出正互信息需求，与其上界矛盾（三维篇 §8） \\
信息论：熵链式法则 & 跨尺度记账（动帧熵增量） & 坐标熵率的定义与下界需求（三维篇 §7，四维篇 §5） \\
分形几何：径向投影 bootstrap & 跨帧参数刚性 & 乘积斜率 $\Delta t\,\Delta\vartheta$ 的钉住估计（三维篇 §9；手稿自述 adapts Shmerkin--Wang 与 Orponen--Shmerkin--Wang）\upcite{shmerkinwang2025,orponensw2024} \\
分形几何：Furstenberg 型估计 & 平面填充与平面刚性 & 近共面短块上的密度输入与保边平面刚性（三维篇 §4，四维篇 §7）\upcite{renwang2025,wangwu2024} \\
分形几何：乘法卷积的 Fourier 衰减 & 标量系数的小区间概率 & 平面刚性工具的前置构造（四维篇 Prop 7.6 之前）\upcite{orponend2024} \\
分形几何：投影与熵平均 & 非集中性与降维 & 标量投影定理（四维篇 Theorem 8.2）与条件非集中估计\upcite{hochmanshmerkin2012} \\
实代数几何：半代数描述、量词消去 & 论证的复杂度纪律 & “温和数据”装置（Definition 2.2）与全局一致有界性（Lemma 3.2 的证明）\upcite{bpr2006} \\
实代数几何：Bernstein--Remez 类 & 多项式拟合的导数与端点控制 & 次数恢复（Prop 6.8）：Remez 迭代 + 四方向 Vandermonde 求逆定出三次项\upcite{roytwaryomdin1997} \\
实代数几何：解析分支与延拓 & 代数叶片的路径延拓 & 端点比较球上的一致范数界（Lemma 3.2；本附录 C.1） \\
\bottomrule
\end{tabularx}
\end{table}

\begin{keybox}[可记忆主句]
三套语言各自承担不可替代的职责：信息论把“失配”变成可计算的信息量，分形几何供给投影与刚性输入，实代数几何规定一致性与常数的纪律——借入的不是定理，而是看问题的方式，而每种方式都被重新证明成了定量形式。
\end{keybox}

% 参考文献前强制清空浮动队列，防表格漂入文献区
\FloatBarrier
% ============================================================
%  参考文献（GB/T 7714 顺序编码制）
%  说明：条目顺序按正文叙事（历史 → 现代 → OpenAI 与核验），
%        终稿前做引用闭合校验（见 build.ps1）
% ============================================================
\begin{thebibliography}{99}
\setlength{\itemsep}{1pt}
\raggedright

% ---------- 历史文献（1917–2002） ----------
\bibitem{fujiwara1917}
FUJIWARA M, KAKEYA S. On some problems of maxima and minima for the curve of constant breadth and the in-revolvable curve of the equilateral triangle[J]. T\^ohoku Mathematical Journal, First Series, 1917, 11: 92--110.

\bibitem{besicovitch1928}
BESICOVITCH A S. On Kakeya's problem and a similar one[J]. Mathematische Zeitschrift, 1928, 27(1): 312--320.

\bibitem{davies1971}
DAVIES R O. Some remarks on the Kakeya problem[J]. Proceedings of the Cambridge Philosophical Society, 1971, 69(3): 417--421.

\bibitem{fefferman1971}
FEFFERMAN C. The multiplier problem for the ball[J]. Annals of Mathematics, 1971, 94(2): 330--336.

\bibitem{cordoba1977}
C\'ORDOBA A. The Kakeya maximal function and the spherical summation multipliers[J]. American Journal of Mathematics, 1977, 99(1): 1--22.

\bibitem{bourgain1991}
BOURGAIN J. Besicovitch type maximal operators and applications to Fourier analysis[J]. Geometric and Functional Analysis, 1991, 1(2): 147--187.

\bibitem{bourgain1999}
BOURGAIN J. On the dimension of Kakeya sets and related maximal inequalities[J]. Geometric and Functional Analysis, 1999, 9(2): 256--282.

\bibitem{wolff1995}
WOLFF T H. An improved bound for Kakeya type maximal functions[J]. Revista Matem\'atica Iberoamericana, 1995, 11(3): 651--674.

\bibitem{wolffbook2003}
WOLFF T H. Lectures on harmonic analysis[M]. Providence: American Mathematical Society, 2003.

\bibitem{klt2000}
KATZ N H, \L ABA I, TAO T. An improved bound on the Minkowski dimension of Besicovitch sets in $\mathbb{R}^3$[J]. Annals of Mathematics, 2000, 152(2): 383--446.

\bibitem{labatao2001}
\L ABA I, TAO T. An improved bound for the Minkowski dimension of Besicovitch sets in medium dimension[J]. Geometric and Functional Analysis, 2001, 11(4): 773--806.

\bibitem{katztai2002}
KATZ N H, TAO T. New bounds for Kakeya problems[J]. Journal d'Analyse Math\'ematique, 2002, 87: 231--263.

% ---------- 多元/多线性与多项式方法（2006–2018） ----------
\bibitem{bct2006}
BENNETT J, CARBERY A, TAO T. On the multilinear restriction and Kakeya conjectures[J]. Acta Mathematica, 2006, 196(2): 261--302.

\bibitem{dvir2009}
DVIR Z. On the size of Kakeya sets in finite fields[J]. Journal of the American Mathematical Society, 2009, 22(4): 1093--1097.

\bibitem{guth2010}
GUTH L. The endpoint case of the Bennett--Carbery--Tao multilinear Kakeya conjecture[J]. Acta Mathematica, 2010, 205(2): 263--286.

\bibitem{bourgainguth2011}
BOURGAIN J, GUTH L. Bounds on oscillatory integral operators based on multilinear estimates[J]. Geometric and Functional Analysis, 2011, 21(6): 1239--1295.

\bibitem{carberyvaldimarsson2013}
CARBERY A, VALDIMARSSON S I. The endpoint multilinear Kakeya theorem via the Borsuk--Ulam theorem[J]. Journal of Functional Analysis, 2013, 264(7): 1643--1663.

\bibitem{guthkatz2015}
GUTH L, KATZ N H. On the Erd\H{o}s distinct distances problem in the plane[J]. Annals of Mathematics, 2015, 181(1): 155--190.

\bibitem{guthzahl2018}
GUTH L, ZAHL J. Polynomial Wolff axioms and Kakeya-type estimates in $\mathbb{R}^4$[J]. Proceedings of the London Mathematical Society, 2018, 117(1): 192--220.

\bibitem{katzrogers2018}
KATZ N H, ROGERS K M. On the polynomial Wolff axioms[J]. Geometric and Functional Analysis, 2018, 28(6): 1706--1716.

\bibitem{zahl2018}
ZAHL J. A discretized Severi-type theorem with applications to harmonic analysis[J]. Geometric and Functional Analysis, 2018, 28(4): 1131--1181.

% ---------- 四维进展与粘性纲领（2019–2026） ----------
\bibitem{katz2019}
KATZ N H, ZAHL J. An improved bound on the Hausdorff dimension of Besicovitch sets in $\mathbb{R}^3$[J]. Journal of the American Mathematical Society, 2019, 32(1): 195--259.

\bibitem{katz2021}
KATZ N H, ZAHL J. A Kakeya maximal function estimate in four dimensions using planebrushes[J]. Revista Matem\'atica Iberoamericana, 2021, 37(1): 317--359.

\bibitem{wangz2022}
WANG H, ZAHL J. Sticky Kakeya sets and the sticky Kakeya conjecture[J]. Journal of the American Mathematical Society, 2026, 39(2): 515--585.

\bibitem{wangz2025assouad}
WANG H, ZAHL J. The Assouad dimension of Kakeya sets in $\mathbb{R}^3$[J]. Inventiones Mathematicae, 2025, 241(1): 153--206.

\bibitem{wangz2025}
WANG H, ZAHL J. Volume estimates for unions of convex sets, and the Kakeya set conjecture in three dimensions[EB/OL]. arXiv:2502.17655, 2025.

\bibitem{gwz2026}
GUTH L, WANG H, ZAHL J. A streamlined proof of the Kakeya set conjecture in $\mathbb{R}^3$[EB/OL]. arXiv:2601.14411, 2026.

\bibitem{renwang2025}
REN K, WANG H. Furstenberg sets estimate in the plane[EB/OL]. arXiv:2308.08819, 2025.

\bibitem{wangwu2024}
WANG H, WU S. Restriction estimates using decoupling theorems and two-ends Furstenberg inequalities[EB/OL]. arXiv:2411.08871, 2024.

\bibitem{borges2025}
BORGES T, CHAN T, CHEN M, LIU D, XI Y, ZHAN Y. Restriction and Kakeya maximal estimates in $\mathbb{R}^4$[EB/OL]. arXiv:2511.22824, 2025.

\bibitem{raichoudhuri2026}
RAI CHOUDHURI M. An improved bound on the Hausdorff dimension of sticky Kakeya sets in $\mathbb{R}^4$[J]. Revista Matem\'atica Iberoamericana, 2026, 42(4): 1227--1256.

\bibitem{wangzakharov2026}
WANG H, ZAKHAROV D. Kakeya and Furstenberg problems for sticky sets of tubes[EB/OL]. arXiv:2609.22035, 2026.

\bibitem{nadjimzadah2026}
NADJIMZADAH A. Bourgain's condition, sticky Kakeya, and new examples[J]. Geometric and Functional Analysis, 2026, 36: 1131--1173.

\bibitem{nadjimzadah2026b}
NADJIMZADAH A. Lifting curved Kakeya sets to linear Kakeya sets[EB/OL]. arXiv:2609.06299, 2026.

\bibitem{gaoliuxi2025}
GAO C, LIU D, XI Y. Curved Kakeya sets and Nikodym problems on manifolds[EB/OL]. arXiv:2503.11574, 2025.

\bibitem{keletimathe2023}
KELETI T, M\'ATH\'E A. Equivalences between different forms of the Kakeya conjecture and duality of Hausdorff and packing dimensions for additive complements[EB/OL]. arXiv:2203.15731, 2023.

\bibitem{zahl2025survey}
ZAHL J. A survey of the Kakeya conjecture, 2000--2025[EB/OL]. arXiv:2512.09397, 2025.

% ---------- 跨门派工具来源 ----------
\bibitem{hochmanshmerkin2012}
HOCHMAN M, SHMERKIN P. Local entropy averages and projections of fractal measures[J]. Annals of Mathematics, 2012, 175(3): 1001--1059.

\bibitem{orponend2024}
ORPONEN T, DE SAXC\'E N, SHMERKIN P. On the Fourier decay of multiplicative convolutions[EB/OL]. arXiv:2309.03068, 2024.

\bibitem{orponensw2024}
ORPONEN T, SHMERKIN P, WANG H. Kaufman and Falconer estimates for radial projections and a continuum version of Beck's theorem[J]. Geometric and Functional Analysis, 2024, 34(1): 164--201.

\bibitem{shmerkinwang2025}
SHMERKIN P, WANG H. On the distance sets spanned by sets of dimension $d/2$ in $\mathbb{R}^d$[J]. Geometric and Functional Analysis, 2025, 35: 283--358.

\bibitem{ssv2005}
SUDAKOV B, SZEMER\'EDI E, VU V H. On a question of Erd\H{o}s and Moser[J]. Duke Mathematical Journal, 2005, 129(1): 129--155.

\bibitem{taovu2006}
TAO T, VU V H. Additive combinatorics[M]. Cambridge: Cambridge University Press, 2006.

\bibitem{bpr2006}
BASU S, POLLACK R, ROY M-F. Algorithms in real algebraic geometry[M]. 2nd ed. Berlin: Springer, 2006.

\bibitem{roytwaryomdin1997}
ROYTWARF N, YOMDIN Y. Bernstein classes[J]. Annales de l'Institut Fourier, 1997, 47(3): 825--858.

\bibitem{coverthomas2006}
COVER T M, THOMAS J A. Elements of information theory[M]. 2nd ed. Hoboken: Wiley, 2006.

% ---------- OpenAI 稿件与发布材料 ----------
\bibitem{openai3d2026}
OpenAI. The Kakeya maximal conjecture in three dimensions[EB/OL]. OpenAI Math Release preprint, OAI:The-Kakeya-maximal-conjecture-in-three-dimensions-September-23-2026, 2026-09-23[2026-10-12]. \url{https://github.com/openai/math}.

\bibitem{openai4d2026}
OpenAI. Every four-dimensional Kakeya set has full Hausdorff dimension[EB/OL]. OpenAI Math Release preprint, OAI:Every-four-dimensional-Kakeya-set-has-full-Hausdorff-dimension-September-24-2026, 2026-09-24[2026-10-12]. \url{https://github.com/openai/math}.

\bibitem{openaielliptic2026}
OpenAI. Elliptic capacity propagation and Fourier restriction to the sphere[EB/OL]. OpenAI Math Release preprint, 2026-09-24[2026-10-12]. \url{https://github.com/openai/math}.

\bibitem{openairelease2026}
OpenAI. Sharing AI progress in mathematics[EB/OL]. OpenAI, 2026-10-06[2026-10-12]. \url{https://openai.com/index/sharing-ai-progress-in-mathematics/}.

\bibitem{openairepo2026}
OpenAI. openai/math: manuscript catalogue, Lean formalization catalogue and release history[EB/OL]. GitHub repository, 2026[2026-10-12]. \url{https://github.com/openai/math}.

\bibitem{openaiten2026}
OpenAI. Ten advances in mathematics and theoretical computer science[EB/OL]. OpenAI, 2026-08-01[2026-10-12]. \url{https://openai.com/index/ten-advances-in-mathematics/}.

\bibitem{openaiunit2026}
OpenAI. An OpenAI model has disproved a central conjecture in discrete geometry[EB/OL]. OpenAI, 2026-05[2026-10-12]. \url{https://openai.com/index/model-disproves-discrete-geometry-conjecture/}.

% ---------- 核验、形式化与学界的反应 ----------
\bibitem{agmai2026}
Advisory Group on Mathematics and Artificial Intelligence, Institute for Advanced Study. Public recommendations[EB/OL]. 2026-09-29[2026-10-12]. \url{https://agmai.org/general-sep29/}.

\bibitem{ahm2026}
Association for Human Mathematics. AHM statement on OpenAI's October 6 release of mathematical documents[EB/OL]. 2026-10-07[2026-10-12]. \url{https://www.ahmath.org/statements}.

\bibitem{tao2026}
TAO T. AHM statement on OpenAI's October 6 release of mathematical documents（转载）[EB/OL]. What's new, 2026-10-07[2026-10-12]. \url{https://terrytao.wordpress.com/2026/10/07/}.

\bibitem{bastounis2026}
BASTOUNIS A, CIRCELLI F, HANSEN A C. Navier--Stokes lost in translation: why Lean verification of AI autoformalisation does not guarantee correct natural language proofs[EB/OL]. arXiv:2610.08144, 2026.

\bibitem{numina2026}
中国科学院数学与系统科学研究院, Project Numina. 三维挂谷猜想证明完成形式化验证[EB/OL]. 2026-09[2026-10-12]. \url{http://www.amss.cas.cn/kyjz1/202609/t20260908_8279602.html}.

\bibitem{xinhua2026}
新华社. 中国数学家邓煜、王虹获得菲尔兹奖[EB/OL]. 2026-07-24[2026-10-12]. \url{https://www.news.cn/world/20260724/4f58ff1485a34463a183f86d33e1ce93/c.html}.

\bibitem{pku2026}
北京大学新闻网. 北大校友王虹、邓煜荣获菲尔兹奖[EB/OL]. 2026-07[2026-10-12]. \url{https://news.pku.edu.cn/xwzh/408c5f40e51b452ab715151c5d0e29b9.htm}.

\bibitem{guangming2026}
光明网. 一根针，能藏进多小的空间——三维挂谷猜想与王虹的数学突破[EB/OL]. 2026-08-27[2026-10-12]. \url{https://news.gmw.cn/2026-08/27/content_38966297.htm}.

\end{thebibliography}

\end{document}